% Options for packages loaded elsewhere
\PassOptionsToPackage{unicode}{hyperref}
\PassOptionsToPackage{hyphens}{url}
\documentclass[
]{article}
\usepackage{xcolor}
\usepackage{amsmath,amssymb}
\setcounter{secnumdepth}{-\maxdimen} % remove section numbering
\usepackage{iftex}
\ifPDFTeX
  \usepackage[T1]{fontenc}
  \usepackage[utf8]{inputenc}
  \usepackage{textcomp} % provide euro and other symbols
\else % if luatex or xetex
  \usepackage{unicode-math} % this also loads fontspec
  \defaultfontfeatures{Scale=MatchLowercase}
  \defaultfontfeatures[\rmfamily]{Ligatures=TeX,Scale=1}
\fi
\usepackage{lmodern}
\ifPDFTeX\else
  % xetex/luatex font selection
\fi
% Use upquote if available, for straight quotes in verbatim environments
\IfFileExists{upquote.sty}{\usepackage{upquote}}{}
\IfFileExists{microtype.sty}{% use microtype if available
  \usepackage[]{microtype}
  \UseMicrotypeSet[protrusion]{basicmath} % disable protrusion for tt fonts
}{}
\makeatletter
\@ifundefined{KOMAClassName}{% if non-KOMA class
  \IfFileExists{parskip.sty}{%
    \usepackage{parskip}
  }{% else
    \setlength{\parindent}{0pt}
    \setlength{\parskip}{6pt plus 2pt minus 1pt}}
}{% if KOMA class
  \KOMAoptions{parskip=half}}
\makeatother
\usepackage{color}
\usepackage{fancyvrb}
\newcommand{\VerbBar}{|}
\newcommand{\VERB}{\Verb[commandchars=\\\{\}]}
\DefineVerbatimEnvironment{Highlighting}{Verbatim}{commandchars=\\\{\}}
% Add ',fontsize=\small' for more characters per line
\newenvironment{Shaded}{}{}
\newcommand{\AlertTok}[1]{\textcolor[rgb]{1.00,0.00,0.00}{\textbf{#1}}}
\newcommand{\AnnotationTok}[1]{\textcolor[rgb]{0.38,0.63,0.69}{\textbf{\textit{#1}}}}
\newcommand{\AttributeTok}[1]{\textcolor[rgb]{0.49,0.56,0.16}{#1}}
\newcommand{\BaseNTok}[1]{\textcolor[rgb]{0.25,0.63,0.44}{#1}}
\newcommand{\BuiltInTok}[1]{\textcolor[rgb]{0.00,0.50,0.00}{#1}}
\newcommand{\CharTok}[1]{\textcolor[rgb]{0.25,0.44,0.63}{#1}}
\newcommand{\CommentTok}[1]{\textcolor[rgb]{0.38,0.63,0.69}{\textit{#1}}}
\newcommand{\CommentVarTok}[1]{\textcolor[rgb]{0.38,0.63,0.69}{\textbf{\textit{#1}}}}
\newcommand{\ConstantTok}[1]{\textcolor[rgb]{0.53,0.00,0.00}{#1}}
\newcommand{\ControlFlowTok}[1]{\textcolor[rgb]{0.00,0.44,0.13}{\textbf{#1}}}
\newcommand{\DataTypeTok}[1]{\textcolor[rgb]{0.56,0.13,0.00}{#1}}
\newcommand{\DecValTok}[1]{\textcolor[rgb]{0.25,0.63,0.44}{#1}}
\newcommand{\DocumentationTok}[1]{\textcolor[rgb]{0.73,0.13,0.13}{\textit{#1}}}
\newcommand{\ErrorTok}[1]{\textcolor[rgb]{1.00,0.00,0.00}{\textbf{#1}}}
\newcommand{\ExtensionTok}[1]{#1}
\newcommand{\FloatTok}[1]{\textcolor[rgb]{0.25,0.63,0.44}{#1}}
\newcommand{\FunctionTok}[1]{\textcolor[rgb]{0.02,0.16,0.49}{#1}}
\newcommand{\ImportTok}[1]{\textcolor[rgb]{0.00,0.50,0.00}{\textbf{#1}}}
\newcommand{\InformationTok}[1]{\textcolor[rgb]{0.38,0.63,0.69}{\textbf{\textit{#1}}}}
\newcommand{\KeywordTok}[1]{\textcolor[rgb]{0.00,0.44,0.13}{\textbf{#1}}}
\newcommand{\NormalTok}[1]{#1}
\newcommand{\OperatorTok}[1]{\textcolor[rgb]{0.40,0.40,0.40}{#1}}
\newcommand{\OtherTok}[1]{\textcolor[rgb]{0.00,0.44,0.13}{#1}}
\newcommand{\PreprocessorTok}[1]{\textcolor[rgb]{0.74,0.48,0.00}{#1}}
\newcommand{\RegionMarkerTok}[1]{#1}
\newcommand{\SpecialCharTok}[1]{\textcolor[rgb]{0.25,0.44,0.63}{#1}}
\newcommand{\SpecialStringTok}[1]{\textcolor[rgb]{0.73,0.40,0.53}{#1}}
\newcommand{\StringTok}[1]{\textcolor[rgb]{0.25,0.44,0.63}{#1}}
\newcommand{\VariableTok}[1]{\textcolor[rgb]{0.10,0.09,0.49}{#1}}
\newcommand{\VerbatimStringTok}[1]{\textcolor[rgb]{0.25,0.44,0.63}{#1}}
\newcommand{\WarningTok}[1]{\textcolor[rgb]{0.38,0.63,0.69}{\textbf{\textit{#1}}}}
\usepackage{longtable,booktabs,array}
\newcounter{none} % for unnumbered tables
\usepackage{calc} % for calculating minipage widths
% Correct order of tables after \paragraph or \subparagraph
\usepackage{etoolbox}
\makeatletter
\patchcmd\longtable{\par}{\if@noskipsec\mbox{}\fi\par}{}{}
\makeatother
% Allow footnotes in longtable head/foot
\IfFileExists{footnotehyper.sty}{\usepackage{footnotehyper}}{\usepackage{footnote}}
\makesavenoteenv{longtable}
\setlength{\emergencystretch}{3em} % prevent overfull lines
\providecommand{\tightlist}{%
  \setlength{\itemsep}{0pt}\setlength{\parskip}{0pt}}
\usepackage{bookmark}
\IfFileExists{xurl.sty}{\usepackage{xurl}}{} % add URL line breaks if available
\urlstyle{same}
\hypersetup{
  hidelinks,
  pdfcreator={LaTeX via pandoc}}

\author{}
\date{}

\begin{document}

\section{Dryad: Distributed Data-Parallel Programs from Sequential
Building
Blocks}\label{dryad-distributed-data-parallel-programs-from-sequential-building-blocks}

Michael Isard

Microsoft Research, Silicon Valley

Mihai Budiu

Microsoft Research, Silicon Valley

Yuan Yu

Microsoft Research, Silicon Valley

Andrew Birrell

Microsoft Research, Silicon Valley

Dennis Fetterly

Microsoft Research, Silicon Valley

\section{ABSTRACT}\label{abstract}

Dryad is a general-purpose distributed execution engine for coarse-grain
data-parallel applications. A Dryad application combines computational
``vertices'' with communication ``channels'' to form a dataflow graph.
Dryad runs the application by executing the vertices of this graph on a
set of available computers, communicating as appropriate through files,
TCP pipes, and shared-memory FIFOs.

The vertices provided by the application developer are quite simple and
are usually written as sequential programs with no thread creation or
locking. Concurrency arises from Dryad scheduling vertices to run
simultaneously on multiple computers, or on multiple CPU cores within a
computer. The application can discover the size and placement of data at
run time, and modify the graph as the computation progresses to make
efficient use of the available resources.

Dryad is designed to scale from powerful multi-core single computers,
through small clusters of computers, to data centers with thousands of
computers. The Dryad execution engine handles all the difficult problems
of creating a large distributed, concurrent application: scheduling the
use of computers and their CPUs, recovering from communication or
computer failures, and transporting data between vertices.

\section{Categories and Subject
Descriptors}\label{categories-and-subject-descriptors}

D.1.3 {[}PROGRAMMING TECHNIQUES{]}: Concurrent Programming---Distributed
programming

\section{General Terms}\label{general-terms}

Performance, Design, Reliability

\section{Keywords}\label{keywords}

Concurrency, Distributed Programming, Dataflow, Cluster Computing

Permission to make digital or hard copies of all or part of this work
for personal or classroom use is granted without fee provided that
copies are not made or distributed for profit or commercial advantage
and that copies bear this notice and the full citation on the first
page. To copy otherwise, to republish, to post on servers or to
redistribute to lists, requires prior specific permission and/or a fee.

EuroSys'07, March 21--23, 2007, Lisboa, Portugal. Copyright 2007 ACM
978-1-59593-636-3/07/0003 \ldots\$5.00.

\section{1. INTRODUCTION}\label{introduction}

The Dryad project addresses a long-standing problem: how can we make it
easier for developers to write efficient parallel and distributed
applications? We are motivated both by the emergence of large-scale
internet services that depend on clusters of hundreds or thousands of
general-purpose servers, and also by the prediction that future advances
in local computing power will come from increasing the number of cores
on a chip rather than improving the speed or instruction-level
parallelism of a single core {[}3{]}. Both of these scenarios involve
resources that are in a single administrative domain, connected using a
known, high-performance communication topology, under centralized
management and control. In such cases many of the hard problems that
arise in wide-area distributed systems may be sidestepped: these include
high-latency and unreliable networks, control of resources by separate
federated or competing entities, and issues of identity for
authentication and access control. Our primary focus is instead on the
simplicity of the programming model and the reliability, efficiency and
scalability of the applications.

For many resource-intensive applications, the simplest way to achieve
scalable performance is to exploit data parallelism. There has
historically been a great deal of work in the parallel computing
community both on systems that automatically discover and exploit
parallelism in sequential programs, and on those that require the
developer to explicitly expose the data dependencies of a computation.
There are still limitations to the power of fully-automatic
parallelization, and so we build mainly on ideas from the latter
research tradition. Condor {[}37{]} was an early example of such a
system in a distributed setting, and we take more direct inspiration
from three other models: shader languages developed for graphic
processing units (GPUs) {[}30, 36{]}, Google's MapReduce system
{[}16{]}, and parallel databases {[}18{]}. In all these programming
paradigms, the system dictates a communication graph, but makes it
simple for the developer to supply subroutines to be executed at
specified graph vertices. All three have demonstrated great success, in
that large numbers of developers have been able to write concurrent
software that is reliably executed in a distributed fashion.

We believe that a major reason for the success of GPU shader languages,
MapReduce and parallel databases is that the developer is explicitly
forced to consider the data parallelism of the computation. Once an
application is cast into this framework, the system is automatically
able to provide the necessary scheduling and distribution. The developer

need have no understanding of standard concurrency mechanisms such as
threads and fine-grain concurrency control, which are known to be
difficult to program correctly. Instead the system runtime abstracts
these issues from the developer, and also deals with many of the hardest
distributed computing problems, most notably resource allocation,
scheduling, and the transient or permanent failure of a subset of
components in the system. By fixing the boundary between the
communication graph and the subroutines that inhabit its vertices, the
model guides the developer towards an appropriate level of granularity.
The system need not try too hard to extract parallelism within a
developer-provided subroutine, while it can exploit the fact that
dependencies are all explicitly encoded in the flow graph to efficiently
distribute the execution across those subroutines. Finally, developers
now work at a suitable level of abstraction for writing scalable
applications since the resources available at execution time are not
generally known at the time the code is written.

The aforementioned systems restrict an application's communication flow
for different reasons. GPU shader languages are strongly tied to an
efficient underlying hardware implementation that has been tuned to give
good performance for common graphics memory-access patterns. MapReduce
was designed to be accessible to the widest possible class of
developers, and therefore aims for simplicity at the expense of
generality and performance. Parallel databases were designed for
relational algebra manipulations (e.g.~SQL) where the communication
graph is implicit.

By contrast, the Dryad system allows the developer fine control over the
communication graph as well as the subroutines that live at its
vertices. A Dryad application developer can specify an arbitrary
directed acyclic graph to describe the application's communication
patterns, and express the data transport mechanisms (files, TCP pipes,
and shared-memory FIFOs) between the computation vertices. This direct
specification of the graph also gives the developer greater flexibility
to easily compose basic common operations, leading to a distributed
analogue of ``piping'' together traditional Unix utilities such as grep,
sort and head.

Dryad is notable for allowing graph vertices (and computations in
general) to use an arbitrary number of inputs and outputs. MapReduce
restricts all computations to take a single input set and generate a
single output set. SQL and shader languages allow multiple inputs but
generate a single output from the user's perspective, though SQL query
plans internally use multiple-output vertices.

In this paper, we demonstrate that careful choices in graph construction
and refinement can substantially improve application performance, while
compromising little on the programmability of the system. Nevertheless,
Dryad is certainly a lower-level programming model than SQL or DirectX.
In order to get the best performance from a native Dryad application,
the developer must understand the structure of the computation and the
organization and properties of the system resources. Dryad was however
designed to be a suitable infrastructure on which to layer simpler,
higher-level programming models. It has already been used, by ourselves
and others, as a platform for several domain-specific systems that are
briefly sketched in Section 7. These rely on Dryad to manage the
complexities of distribution, scheduling, and fault-tolerance, but hide
many of the details of the underlying system from the application
developer. They use heuristics to automatically select and tune
appropriate Dryad features, and thereby get good performance for most
simple applications.

We summarize Dryad's contributions as follows:

- We built a general-purpose, high performance distributed execution
engine. The Dryad execution engine handles many of the difficult
problems of creating a large distributed, concurrent application:
scheduling across resources, optimizing the level of concurrency within
a computer, recovering from communication or computer failures, and
delivering data to where it is needed. Dryad supports multiple different
data transport mechanisms between computation vertices and explicit
dataflow graph construction and refinement.

- We demonstrated the excellent performance of Dryad from a single
multi-core computer up to clusters consisting of thousands of computers
on several nontrivial, real examples. We further demonstrated that
Dryad's fine control over an application's dataflow graph gives the
programmer the necessary tools to optimize trade-offs between
parallelism and data distribution overhead. This validated Dryad's
design choices.

- We explored the programmability of Dryad on two fronts. First, we have
designed a simple graph description language that empowers the developer
with explicit graph construction and refinement to fully take advantage
of the rich features of the Dryad execution engine. Our user experiences
lead us to believe that, while it requires some effort to learn, a
programmer can master the APIs required for most of the applications in
a couple of weeks. Second, we (and others within Microsoft) have built
simpler, higher-level programming abstractions for specific application
domains on top of Dryad. This has significantly lowered the barrier to
entry and increased the acceptance of Dryad among domain experts who are
interested in using Dryad for rapid application prototyping. This
further validated Dryad's design choices.

The next three sections describe the abstract form of a Dryad
application and outline the steps involved in writing one. The Dryad
scheduler is described in Section 5; it handles all of the work of
deciding which physical resources to schedule work on, routing data
between computations, and automatically reacting to computer and network
failures. Section 6 reports on our experimental evaluation of the
system, showing its flexibility and scaling characteristics in a small
cluster of 10 computers, as well as details of larger-scale experiments
performed on clusters with thousands of computers. We conclude in
Sections 8 and 9 with a discussion of the related literature and of
future research directions.

\section{2. SYSTEM OVERVIEW}\label{system-overview}

The overall structure of a Dryad job is determined by its communication
flow. A job is a directed acyclic graph where each vertex is a program
and edges represent data channels. It is a logical computation graph
that is automatically mapped onto physical resources by the runtime. In
particular, there may be many more vertices in the graph than execution
cores in the computing cluster.

At run time each channel is used to transport a finite sequence of
structured items. This channel abstraction has several concrete
implementations that use shared memory, TCP pipes, or files temporarily
persisted in a file system. As far as the program in each vertex is
concerned, channels produce and consume heap objects that inherit from a
base type. This means that a vertex program reads and writes its data in
the same way regardless of whether a channel serializes its data to
buffers on a disk or TCP stream, or passes object pointers directly via
shared memory. The Dryad system does not include any native data model
for serialization and the concrete type of an item is left entirely up
to applications, which can supply their own serialization and
deserialization routines. This decision allows us to support
applications that operate directly on existing data including exported
SQL tables and textual log files. In practice most applications use one
of a small set of library item types that we supply such as
newline-terminated text strings and tuples of base types.

A schematic of the Dryad system organization is shown in Figure 1. A
Dryad job is coordinated by a process called the ``job manager''
(denoted JM in the figure) that runs either within the cluster or on a
user's workstation with network access to the cluster. The job manager
contains the application-specific code to construct the job's
communication graph along with library code to schedule the work across
the available resources. All data is sent directly between vertices and
thus the job manager is only responsible for control decisions and is
not a bottleneck for any data transfers.

\textit{[Image omitted: images/e599d59de9c5e598163fe679c6a0039e62b1182af990ccc9df9ccd03442741b4.jpg]}

flowchart

\begin{Shaded}
\begin{Highlighting}[]
\NormalTok{graph TD}
\NormalTok{    A["Job schedule"] {-}{-}\textgreater{} B["JSM"]}
\NormalTok{    B {-}{-}\textgreater{} C["Control plane"]}
\NormalTok{    C {-}{-}\textgreater{} D["Data plane"]}
\NormalTok{    D {-}{-}\textgreater{} E["Files, FIFO, Network"]}
\NormalTok{    E {-}{-}\textgreater{} F["V"]}
\NormalTok{    F {-}{-}\textgreater{} G["D"]}
\NormalTok{    G {-}{-}\textgreater{} H["D"]}
\NormalTok{    H {-}{-}\textgreater{} I["D"]}
\NormalTok{    I {-}{-}\textgreater{} J["V"]}
\NormalTok{    J {-}{-}\textgreater{} K["D"]}
\NormalTok{    K {-}{-}\textgreater{} L["V"]}
\NormalTok{    L {-}{-}\textgreater{} M["Control plane"]}
\NormalTok{    M {-}{-}\textgreater{} N["NS"]}
\end{Highlighting}
\end{Shaded}

Figure 1: The Dryad system organization. The job manager (JM) consults
the name server (NS) to discover the list of available computers. It
maintains the job graph and schedules running vertices (V) as computers
become available using the daemon (D) as a proxy. Vertices exchange data
through files, TCP pipes, or shared-memory channels. The shaded bar
indicates the vertices in the job that are currently running.

The cluster has a name server (NS) that can be used to enumerate all the
available computers. The name server also exposes the position of each
computer within the network topology so that scheduling decisions can
take account of locality. There is a simple daemon (D) running on each
computer in the cluster that is responsible for creating processes on
behalf of the job manager. The first time a vertex (V) is executed on a
computer its binary is sent from the job manager to the daemon and
subsequently it is executed from a cache. The daemon acts as a proxy so
that the job manager can communicate with the remote vertices and
monitor the state of the computation and how much data has been read and
written on its channels. It is straightforward to run a name server and
a set of daemons on a user workstation to simulate a cluster and thus
run an entire job locally while debugging.

A simple task scheduler is used to queue batch jobs. We use a
distributed storage system, not described here, that shares with the
Google File System {[}21{]} the property that large files can be broken
into small pieces that are replicated and distributed across the local
disks of the cluster computers. Dryad also supports the use of NTFS for
accessing files directly on local computers, which can be convenient for
small clusters with low management overhead.

\section{2.1 An example SQL query}\label{an-example-sql-query}

In this section, we describe a concrete example of a Dryad application
that will be further developed throughout the remainder of the paper.
The task we have chosen is representative of a new class of eScience
applications, where scientific investigation is performed by processing
large amounts of data available in digital form {[}24{]}. The database
that we use is derived from the Sloan Digital Sky Survey (SDSS),
available online at http://skyserver.sdss.org.

We chose the most time consuming query (Q18) from a published study
based on this database {[}23{]}. The task is to identify a
``gravitational lens'' effect: it finds all the objects in the database
that have neighboring objects within 30 arc seconds such that at least
one of the neighbors has a color similar to the primary object's color.
The query can be expressed in SQL as:

select distinct p.objID from photoObjAll p join neighbors n --- call
this join ``X'' on p.objID = n.objID and n.objID \textless{}
n.neighborObjID and p.mode = 1 join photoObjAll 1 --- call this join
``Y'' on l.objid = n.neighborObjID and l.mode = 1 and
abs((p.u-p.g)-(1.u-1.g))\textless0.05 and
abs((p.g-p.r)-(1.g-1.r))\textless0.05 and
abs((p.r-p.i)-(1.r-1.i))\textless0.05 and
abs((p.i-p.z)-(1.i-1.z))\textless0.05

There are two tables involved. The first, photoObjAll has 354,254,163
records, one for each identified astronomical object, keyed by a unique
identifier objID. These records also include the object's color, as a
magnitude (logarithmic brightness) in five bands: u, g, r, i and z. The
second table, neighbors has 2,803,165,372 records, one for each object
located within 30 arc seconds of another object. The mode predicates in
the query select only ``primary'' objects. The \textless{} predicate
eliminates duplication caused by the neighbors relationship being
symmetric. The output of joins ``X'' and ``Y'' are 932,820,679 and
83,798 records respectively, and the final hash emits 83,050 records.

The query uses only a few columns from the tables (the complete
photoObjAll table contains 2 KBytes per record). When executed by
SQLServer the query uses an index on photoObjAll keyed by objID with
additional columns for mode, u, g, r, i and z, and an index on neighbors
keyed by objID with an additional neighborObjID column. SQLServer reads
just these indexes, leaving the remainder of the tables' data resting
quietly on disk. (In our experimental setup we in fact omitted unused
columns from the table, to avoid transporting the entire multi-terabyte
database across

the country.) For the equivalent Dryad computation we extracted these
indexes into two binary files, ``ugriz.bin'' and ``neighbors.bin,'' each
sorted in the same order as the indexes. The ``ugriz.bin'' file has
36-byte records, totaling 11.8 GBytes; ``neighbors.bin'' has 16-byte
records, totaling 41.8 GBytes. The output of join ``X'' totals 31.3
GBytes, the output of join ``Y'' is 655 KBytes and the final output is
649 KBytes.

We mapped the query to the Dryad computation shown in Figure 2. Both
data files are partitioned into n approximately equal parts (that we
call \(U_{1}\) through \(U_{n}\) and \(N_{1}\) through \(N_{n}\) ) by
objID ranges, and we use custom C++ item objects for each data record in
the graph. The vertices \(X_{i}\) (for \(1 \leq i \leq n\) ) implement
join ``X'' by taking their partitioned \(U_{i}\) and \(N_{i}\) inputs
and merging them (keyed on objID and filtered by the \textless{}
expression and p.mode=1) to produce records containing objID,
neighborObjID, and the color columns corresponding to objID. The D
vertices distribute their output records to the M vertices, partitioning
by neighborObjID using a range partitioning function four times finer
than that used for the input files. The number four was chosen so that
four pipelines will execute in parallel on each computer, because our
computers have four processors each. The M vertices perform a
non-deterministic merge of their inputs and the S vertices sort on
neighborObjID using an in-memory Quicksort. The output records from
\(S_{4i-3} \ldots S_{4i}\) (for i = 1 through n) are fed into \(Y_i\)
where they are merged with another read of \(U_i\) to implement join
``Y''. This join is keyed on objID (from U) = neighborObjID (from S),
and is filtered by the remainder of the predicate, thus matching the
colors. The outputs of the Y vertices are merged into a hash table at
the H vertex to implement the distinct keyword in the query. Finally, an
enumeration of this hash table delivers the result. Later in the paper
we include more details about the implementation of this Dryad program.

\section{3. DESCRIBING A DRYAD GRAPH}\label{describing-a-dryad-graph}

We have designed a simple language that makes it easy to specify
commonly-occurring communication idioms. It is currently ``embedded''
for convenience in C++ as a library using a mixture of method calls and
operator overloading.

Graphs are constructed by combining simpler subgraphs using a small set
of operations shown in Figure 3. All of the operations preserve the
property that the resulting graph is acyclic. The basic object in the
language is a graph:

\[
G = \langle V _ {G}, E _ {G}, I _ {G}, O _ {G} \rangle .
\]

\textit{[Image omitted: images/4515b9cee1567332b2a95e7f08220bbd572c15ea7de52965c8e3e40447b0bad2.jpg]}

flowchart

Neural network architecture diagram showing connections between input
nodes U, hidden state H, and output nodes X, Y, S, M, D with labeled
distances.

Figure 2: The communication graph for an SQL query. Details are in
Section 2.1.

G contains a sequence of vertices \(V_{G}\) , a set of directed edges
\(E_{G}\) , and two sets \(I_{G} \subseteq V_{G}\) and
\(O_{G} \subseteq V_{G}\) that ``tag'' some of the vertices as being
inputs and outputs respectively. No graph can contain a directed edge
entering an input vertex in \(I_{G}\) , nor one leaving an output vertex
in \(O_{G}\) , and these tags are used below in composition operations.
The input and output edges of a vertex are ordered so an edge connects
specific ``ports'' on a pair of vertices, and a given pair of vertices
may be connected by multiple edges.

\section{3.1 Creating new vertices}\label{creating-new-vertices}

The Dryad libraries define a C++ base class from which all vertex
programs inherit. Each such program has a textual name (which is unique
within an application) and a static ``factory'' that knows how to
construct it. A graph vertex is created by calling the appropriate
static program factory. Any required vertex-specific parameters can be
set at this point by calling methods on the program object. These
parameters are then marshaled along with the unique vertex name to form
a simple closure that can be sent to a remote process for execution.

A singleton graph is generated from a vertex v as
\(G = \langle(v), \emptyset, \{v\}, \{v\} \rangle\) . A graph can be
cloned into a new graph containing k copies of its structure using the
\(\hat{}\) operator where \(C = G^{\hat{}}k\) is defined as:

\[
C = \left\langle V _ {G} ^ {1} \oplus \dots \oplus V _ {G} ^ {k}, E _ {G} ^ {1} \cup \dots \cup E _ {G} ^ {k}, \right.
\]

\[
I _ {G} ^ {1} \cup \dots \cup I _ {G} ^ {k}, O _ {G} ^ {1} \cup \dots \cup O _ {G} ^ {k} \rangle .
\]

Here
\(G^{n} = \langle V_{G}^{n}, E_{G}^{n}, I_{G}^{n}, O_{G}^{n} \rangle\)
is a ``clone'' of G containing copies of all of G's vertices and edges,
\(\oplus\) denotes sequence concatenation, and each cloned vertex
inherits the type and parameters of its corresponding vertex in G.

\section{3.2 Adding graph edges}\label{adding-graph-edges}

New edges are created by applying a composition operation to two
existing graphs. There is a family of compositions all sharing the same
basic structure: C = A ∘ B creates a new graph:

\[
C = \left\langle V _ {A} \oplus V _ {B}, E _ {A} \cup E _ {B} \cup E _ {\text { new }}, I _ {A}, O _ {B} \right\rangle
\]

where C contains the union of all the vertices and edges in A and B,
with A's inputs and B's outputs. In addition, directed edges \(E_{new}\)
are introduced between vertices in \(O_{A}\) and \(I_{B}\) . \(V_{A}\)
and \(V_{B}\) are enforced to be disjoint at run time, and since A and B
are both acyclic, C is also.

Compositions differ in the set of edges \(E_{new}\) that they add into
the graph. We define two standard compositions:

\begin{itemize}
\tightlist
\item
  A \textgreater= B forms a pointwise composition as shown in Figure
  3(c). If \(|O_A| \geq |I_B|\) then a single outgoing edge is created
  from each of \(A\) 's outputs. The edges are assigned in round-robin
  to \(B\) 's inputs. Some of the vertices in \(I_B\) may end up with
  more than one incoming edge. If \(|I_B| > |O_A|\) , a single incoming
  edge is created to each of \(B\) 's inputs, assigned in round-robin
  from \(A\) 's outputs.
\end{itemize}

- A \textgreater\textgreater{} B forms the complete bipartite graph
between \(O_A\) and \(I_B\) and is shown in Figure 3(d).

We allow the user to extend the language by implementing new composition
operations.

\textit{[Image omitted: images/ff989cdf70ec10bc1e1c0ed9215c94118b3632c75e9152c6d45d37a4a104caef.jpg]}

text\_image

a A n A AS = A\^{}n

\textit{[Image omitted: images/5970dbd613120aeafed12b35f45895e1c713b0b495b169f773753f18eb5ad093.jpg]}

flowchart

\begin{Shaded}
\begin{Highlighting}[]
\NormalTok{graph TD}
\NormalTok{    A["A"] {-}{-}\textgreater{} B["B"]}
\NormalTok{    B {-}{-}\textgreater{} A}
\NormalTok{    A {-}{-}\textgreater{}|n| B}
\NormalTok{    B {-}{-}\textgreater{}|n| A}
\NormalTok{    style A fill:\#fff,stroke:\#000}
\NormalTok{    style B fill:\#fff,stroke:\#000}
\NormalTok{    style A fill:\#fff,stroke:\#000}
\NormalTok{    style B fill:\#fff,stroke:\#000}
\NormalTok{    note bottom}
\NormalTok{        AS \textgreater{}= BS}
\NormalTok{    end}
\end{Highlighting}
\end{Shaded}

\textit{[Image omitted: images/dce434e262364650d087d6a56c34d7dabbc07ae4a73bb0ebc8b24fea16bb2fd2.jpg]}

flowchart

Diagram showing relationships between nodes A and B with labeled edges
`n', including a condition `AS \textgreater\textgreater{} BS' below

\textit{[Image omitted: images/478669caa0400c6d242479da1c624231825f0042ccb2bb62e292fceafc228bdd.jpg]}

flowchart

\begin{Shaded}
\begin{Highlighting}[]
\NormalTok{graph TD}
\NormalTok{    B {-}{-}\textgreater{} C}
\NormalTok{    B {-}{-}\textgreater{} D}
\NormalTok{    C {-}{-}\textgreater{} B}
\NormalTok{    D {-}{-}\textgreater{} B}
\NormalTok{    style B fill:\#fff,stroke:\#000}
\NormalTok{    style C fill:\#fff,stroke:\#000}
\NormalTok{    style D fill:\#fff,stroke:\#000}
\NormalTok{    note bottom of B (B \textgreater{}= C) || (B \textgreater{}= D)}
\end{Highlighting}
\end{Shaded}

\textit{[Image omitted: images/d33c64b83adf299ea087e871c5aba18c7b903e81e2092822d22f68e0fb106d39.jpg]}

text\_image

b B n B BS = B\^{}n

\textit{[Image omitted: images/cab46b3cad1a23b9ce3b9db13333f80d1268b2a28a54eb3cb93ba1fa5488ab1d.jpg]}

flowchart

\begin{Shaded}
\begin{Highlighting}[]
\NormalTok{graph TD}
\NormalTok{    A["A"] {-}{-}\textgreater{} C["C"]}
\NormalTok{    B["B"] {-}{-}\textgreater{} C["C"]}
\NormalTok{    A["A"] {-}{-}\textgreater{} C["C"]}
\NormalTok{    B["B"] {-}{-}\textgreater{} C["C"]}
\NormalTok{    A["A"] {-}{-}\textgreater{} C["C"]}
\NormalTok{    B["B"] {-}{-}\textgreater{} C["C"]}
\NormalTok{    A["A"] {-}{-}\textgreater{} C["C"]}
\NormalTok{    B["B"] {-}{-}\textgreater{} C["C"]}
\NormalTok{    C["C"] {-}{-}\textgreater{} A["A"]}
\NormalTok{    C["C"] {-}{-}\textgreater{} B["B"]}
\NormalTok{    C["C"] {-}{-}\textgreater{} A["A"]}
\NormalTok{    style C fill:\#ff0000,stroke:\#333}
\NormalTok{    note right of C E = (AS \textgreater{}= C \textgreater{}= BS)}
\end{Highlighting}
\end{Shaded}

\textit{[Image omitted: images/89aaac0f820a684e01b9a3fe06c825cf681cf4a32aa0f80d05d08757abdafe93.jpg]}

flowchart

\begin{Shaded}
\begin{Highlighting}[]
\NormalTok{graph TD}
\NormalTok{    A["A"] {-}{-}\textgreater{} B["B"]}
\NormalTok{    A {-}{-}\textgreater{} C["C"]}
\NormalTok{    B {-}{-}\textgreater{} C}
\NormalTok{    B {-}{-}\textgreater{} A}
\NormalTok{    C {-}{-}\textgreater{} A}
\NormalTok{    C {-}{-}\textgreater{} B}
\NormalTok{    C {-}{-}\textgreater{} B}
\NormalTok{    A {-}{-}\textgreater{} B}
\NormalTok{    A {-}{-}\textgreater{} C}
\NormalTok{    B {-}{-}\textgreater{} C}
\NormalTok{    style A fill:\#fff,stroke:\#000}
\NormalTok{    style B fill:\#fff,stroke:\#000}
\NormalTok{    style C fill:\#fff,stroke:\#000}
\NormalTok{    style A fill:\#fff,stroke:\#000}
\NormalTok{    style B fill:\#fff,stroke:\#000}
\NormalTok{    style C fill:\#fff,stroke:\#000}
\NormalTok{    style A fill:\#fff,stroke:\#000}
\NormalTok{    style B fill:\#fff,stroke:\#000}
\NormalTok{    style C fill:\#fff,stroke:\#000}
\NormalTok{    style A fill:\#fff,stroke:\#000}
\NormalTok{    style B fill:\#fff,stroke:\#000}
\NormalTok{    style C fill:\#fff,stroke:\#000}
\NormalTok{    style A fill:\#fff,stroke:\#000}
\NormalTok{    style B fill:\#fff,stroke:\#000}
\NormalTok{    style C fill:\#fff,stroke:\#000}
\NormalTok{    style A fill:\#fff,stroke:\#000}
\NormalTok{    style B fill:\#fff,stroke:\#000}
\NormalTok{    style C fill:\#fff,stroke:\#000}
\NormalTok{    style A fill:\#fff,stroke:\#000}
\NormalTok{    style B fill:\#fff,stroke:\#000}
\NormalTok{    style C fill:\#fff,stroke:\#000}
\NormalTok{    style A fill:\#fff,stroke:\#000}
\NormalTok{    style B fill:\#fff,stroke:\#000}
\NormalTok{    style C fill:\#fff,stroke:\#000}
\NormalTok{    style A fill:\#fff,stroke:\#000}
\NormalTok{    style B fill:\#fff,stroke:\#000}
\NormalTok{    style C fill:\#fff,stroke:\#000}
\NormalTok{    style A fill:\#fff,stroke:\#000}
\NormalTok{    style B fill:\#fff,stroke:\#000}
\NormalTok{    style C fill:\#fff,stroke:\#000}
\NormalTok{    style A fill:\#fff,stroke:\#000}
\NormalTok{    style B fill:\#fff,stroke:\#000}
\NormalTok{    style C fill:\#fff,stroke:\#000}
\end{Highlighting}
\end{Shaded}

\textit{[Image omitted: images/6b4eb91acb35349a9f1cb8b09eb79247564c6590b5a5f9bf1f5b0ccec8a1b385.jpg]}

flowchart

\begin{Shaded}
\begin{Highlighting}[]
\NormalTok{graph TD}
\NormalTok{    A {-}{-}\textgreater{} C}
\NormalTok{    C {-}{-}\textgreater{} D}
\NormalTok{    D {-}{-}\textgreater{} B}
\NormalTok{    B {-}{-}\textgreater{} F}
\NormalTok{    F {-}{-}\textgreater{} A}
\NormalTok{    A {-}{-}\textgreater{} A}
\NormalTok{    style A fill:\#fff,stroke:\#000}
\NormalTok{    style B fill:\#fff,stroke:\#000}
\NormalTok{    style C fill:\#fff,stroke:\#000}
\NormalTok{    style D fill:\#fff,stroke:\#000}
\NormalTok{    style F fill:\#fff,stroke:\#000}
\end{Highlighting}
\end{Shaded}

Figure 3: The operators of the graph description language. Circles are
vertices and arrows are graph edges. A triangle at the bottom of a
vertex indicates an input and one at the top indicates an output. Boxes
(a) and (b) demonstrate cloning individual vertices using the \^{}
operator. The two standard connection operations are pointwise
composition using \textgreater= shown in (c) and complete bipartite
composition using \textgreater\textgreater{} shown in (d). (e)
illustrates a merge using \textbar\textbar. The second line of the
figure shows more complex patterns. The merge in (g) makes use of a
``subroutine'' from (f) and demonstrates a bypass operation. For
example, each A vertex might output a summary of its input to C which
aggregates them and forwards the global statistics to every B. Together
the B vertices can then distribute the original dataset (received from
A) into balanced partitions. An asymmetric fork/join is shown in (h).

\section{3.3 Merging two graphs}\label{merging-two-graphs}

The final operation in the language is \textbar\textbar, which merges
two graphs. C = A \textbar\textbar{} B creates a new graph:

\[
C = \left\langle V _ {A} \oplus^ {*} V _ {B}, E _ {A} \cup E _ {B}, I _ {A} \cup^ {*} I _ {B}, O _ {A} \cup^ {*} O _ {B} \right\rangle
\]

where, in contrast to the composition operations, it is not required
that A and B be disjoint. \(V_{A} \oplus^{*} V_{B}\) is the
concatenation of \(V_{A}\) and \(V_{B}\) with duplicates removed from
the second sequence. \(I_{A} \cup^{*} I_{B}\) means the union of A and
B's inputs, minus any vertex that has an incoming edge following the
merge (and similarly for the output case). If a vertex is contained in
\(V_{A} \cap V_{B}\) its input and output edges are concatenated so that
the edges in \(E_{A}\) occur first (with lower port numbers). This
simplification forbids certain graphs with ``crossover'' edges, however
we have not found this restriction to be a problem in practice. The
invariant that the merged graph be acyclic is enforced by a run-time
check.

The merge operation is extremely powerful and makes it easy to construct
typical patterns of communication such as fork/join and bypass as shown
in Figures 3(f)--(h). It also provides the mechanism for assembling a
graph ``by hand'' from a collection of vertices and edges. So for
example, a tree with four vertices a, b, c, and d might be constructed
as \(G = (a \geq b) \mid \mid (b \geq c) \mid \mid (b \geq d)\) .

The graph builder program to construct the query graph in Figure 2 is
shown in Figure 4.

\section{3.4 Channel types}\label{channel-types}

By default each channel is implemented using a temporary file: the
producer writes to disk (typically on its local computer) and the
consumer reads from that file.

In many cases multiple vertices will fit within the resources of a
single computer so it makes sense to execute them all within the same
process. The graph language has an ``encapsulation'' command that takes
a graph G and returns a new vertex \(v_{G}\) . When \(v_{G}\) is run as
a vertex program, the job manager passes it a serialization of G as an
invocation parameter, and it runs all the vertices of G simultaneously
within the same process, connected by edges implemented using
shared-memory FIFOs. While it would always be possible to write a custom
vertex program with the same semantics as G, allowing encapsulation
makes it efficient to combine simple library vertices at the graph layer
rather than re-implementing their functionality as a new vertex program.

Sometimes it is desirable to place two vertices in the same process even
though they cannot be collapsed into a single graph vertex from the
perspective of the scheduler. For example, in Figure 2 the performance
can be improved by placing the first D vertex in the same process as the
first four M and S vertices and thus avoiding some disk I/O, however the
S vertices cannot be started until all of the D vertices complete.

When creating a set of graph edges, the user can optionally specify the
transport protocol to be used. The available protocols are listed in
Table 1. Vertices that are connected using shared-memory channels are
executed within a single process, though they are individually started
as their inputs become available and individually report completion.

Because the dataflow graph is acyclic, scheduling deadlock is impossible
when all channels are either written to temporary files or use
shared-memory FIFOs hidden within

\begin{Shaded}
\begin{Highlighting}[]
\NormalTok{GraphBuilder XSet }\OperatorTok{=}\NormalTok{ moduleX}\OperatorTok{\^{}}\NormalTok{N}\OperatorTok{;}
\NormalTok{GraphBuilder DSet }\OperatorTok{=}\NormalTok{ moduleD}\OperatorTok{\^{}}\NormalTok{N}\OperatorTok{;}
\NormalTok{GraphBuilder MSet }\OperatorTok{=}\NormalTok{ moduleM}\OperatorTok{\^{}(}\NormalTok{N}\OperatorTok{*}\DecValTok{4}\OperatorTok{);}
\NormalTok{GraphBuilder SSet }\OperatorTok{=}\NormalTok{ moduleS}\OperatorTok{\^{}(}\NormalTok{N}\OperatorTok{*}\DecValTok{4}\OperatorTok{);}
\NormalTok{GraphBuilder YSet }\OperatorTok{=}\NormalTok{ moduleY}\OperatorTok{\^{}}\NormalTok{N}\OperatorTok{;}
\NormalTok{GraphBuilder HSet }\OperatorTok{=}\NormalTok{ moduleH}\OperatorTok{\^{}}\DecValTok{1}\OperatorTok{;}

\NormalTok{GraphBuilder XInputs }\OperatorTok{=} \OperatorTok{(}\NormalTok{ugriz1 }\OperatorTok{\textgreater{}=}\NormalTok{ XSet}\OperatorTok{)} \OperatorTok{||} \OperatorTok{(}\NormalTok{neighbor }\OperatorTok{\textgreater{}=}\NormalTok{ XSet}\OperatorTok{);}
\NormalTok{GraphBuilder YInputs }\OperatorTok{=}\NormalTok{ ugriz2 }\OperatorTok{\textgreater{}=}\NormalTok{ YSet}\OperatorTok{;}

\NormalTok{GraphBuilder XToY }\OperatorTok{=}\NormalTok{ XSet }\OperatorTok{\textgreater{}=}\NormalTok{ DSet }\OperatorTok{\textgreater{}\textgreater{}}\NormalTok{ MSet }\OperatorTok{\textgreater{}=}\NormalTok{ SSet}\OperatorTok{;}
\ControlFlowTok{for} \OperatorTok{(}\NormalTok{i }\OperatorTok{=} \DecValTok{0}\OperatorTok{;}\NormalTok{ i }\OperatorTok{\textless{}}\NormalTok{ N}\OperatorTok{*}\DecValTok{4}\OperatorTok{;} \OperatorTok{++}\NormalTok{i}\OperatorTok{)}
\OperatorTok{\{}
\NormalTok{    XToY }\OperatorTok{=}\NormalTok{ XToY }\OperatorTok{||} \OperatorTok{(}\NormalTok{SSet}\OperatorTok{.}\NormalTok{GetVertex}\OperatorTok{(}\NormalTok{i}\OperatorTok{)} \OperatorTok{\textgreater{}=}\NormalTok{ YSet}\OperatorTok{.}\NormalTok{GetVertex}\OperatorTok{(}\NormalTok{i}\OperatorTok{/}\DecValTok{4}\OperatorTok{));}
\OperatorTok{\}}

\NormalTok{GraphBuilder YToH }\OperatorTok{=}\NormalTok{ YSet }\OperatorTok{\textgreater{}=}\NormalTok{ HSet}\OperatorTok{;}
\NormalTok{GraphBuilder HOutputs }\OperatorTok{=}\NormalTok{ HSet }\OperatorTok{\textgreater{}=}\NormalTok{ output}\OperatorTok{;}

\NormalTok{GraphBuilder final }\OperatorTok{=}\NormalTok{ XInputs }\OperatorTok{||}\NormalTok{ YInputs }\OperatorTok{||}\NormalTok{ XToY }\OperatorTok{||}\NormalTok{ YToH }\OperatorTok{||}\NormalTok{ HOutputs}\OperatorTok{;} 
\end{Highlighting}
\end{Shaded}

Figure 4: An example graph builder program. The communication graph
generated by this program is shown in Figure 2.

Channel protocol

Discussion

File (the default)

Preserved after vertex execution until the job completes.

TCP pipe

Requires no disk accesses, but both end-point vertices must be scheduled
to run at the same time.

Shared-memory FIFO

Extremely low communication cost, but end-point vertices must run within
the same process.

Table 1: Channel types.

encapsulated acyclic subgraphs. However, allowing the developer to use
pipes and ``visible'' FIFOs can cause deadlocks. Any connected component
of vertices communicating using pipes or FIFOs must all be scheduled in
processes that are concurrently executing, but this becomes impossible
if the system runs out of available computers in the cluster.

This breaks the abstraction that the user need not know the physical
resources of the system when writing the application. We believe that it
is a worthwhile trade-off, since, as reported in our experiments in
Section 6, the resulting performance gains can be substantial. Note also
that the system could always avoid deadlock by ``downgrading'' a pipe
channel to a temporary file, at the expense of introducing an unexpected
performance cliff.

\section{3.5 Job inputs and outputs}\label{job-inputs-and-outputs}

Large input files are typically partitioned and distributed across the
computers of the cluster. It is therefore natural to group a logical
input into a graph
\(G = \langle V_{P}, \emptyset, \emptyset, V_{P} \rangle\) where
\(V_{P}\) is a sequence of ``virtual'' vertices corresponding to the
partitions of the input. Similarly on job completion a set of output
partitions can be logically concatenated to form a single named
distributed file. An application will generally interrogate its input
graphs to read the number of partitions at run time and automatically
generate the appropriately replicated graph.

\section{3.6 Job Stages}\label{job-stages}

When the graph is constructed every vertex is placed in a ``stage'' to
simplify job management. The stage topology can be seen as a
``skeleton'' or summary of the overall job, and the stage topology of
our example Skyserver query application is shown in Figure 5. Each
distinct type of vertex is grouped into a separate stage. Most stages
are connected using the \textgreater= operator, while D is connected to
M using the \textgreater\textgreater{} operator. The skeleton is used as
a guide for generating summaries when monitoring a job, and can also be
exploited by the automatic optimizations described in Section 5.2.

\section{4. WRITING A VERTEX PROGRAM}\label{writing-a-vertex-program}

The primary APIs for writing a Dryad vertex program are exposed through
C++ base classes and objects. It was a design requirement for Dryad
vertices to be able to incorporate legacy source and libraries, so we
deliberately avoided adopting any Dryad-specific language or sandboxing
restrictions. Most of the existing code that we anticipate integrating
into vertices is written in C++, but it is straightforward to implement
API wrappers so that developers can write vertices in other languages,
for example C\#. There is also significant value for some domains in
being able to run unmodified legacy executables in vertices, and so we
support this as explained in Section 4.2 below.

\section{4.1 Vertex execution}\label{vertex-execution}

Dryad includes a runtime library that is responsible for setting up and
executing vertices as part of a distributed computation. As outlined in
Section 3.1 the runtime receives a closure from the job manager
describing the vertex to be run, and URIs describing the input and
output channels to connect to it. There is currently no type-checking
for channels and the vertex must be able to deter-

mine, either statically or from the invocation parameters, the types of
the items that it is expected to read and write

\textit{[Image omitted: images/1ec0239ecf846265062658e73a42e67da1cfd8d40bf9d12663f438de50ad309a.jpg]}

flowchart

\begin{Shaded}
\begin{Highlighting}[]
\NormalTok{graph TD}
\NormalTok{    H {-}{-}\textgreater{} Y}
\NormalTok{    Y {-}{-}\textgreater{} S}
\NormalTok{    S {-}{-}\textgreater{} M}
\NormalTok{    M {-}{-}\textgreater{} D}
\NormalTok{    D {-}{-}\textgreater{} X}
\NormalTok{    X {-}{-}\textgreater{} U}
\NormalTok{    X {-}{-}\textgreater{} N}
\NormalTok{    U {-}{-}\textgreater{} H}
\NormalTok{    N {-}{-}\textgreater{} X}
\NormalTok{    U {-}{-}\textgreater{} Y}
\NormalTok{    Y {-}{-}\textgreater{} H}
\end{Highlighting}
\end{Shaded}

Figure 5: The stages of the Dryad computation from Figure 2. Section 3.6
has details.

on each channel in order to supply the correct serialization routines.
The body of a vertex is invoked via a standard Main method that includes
channel readers and writers in its argument list. The readers and
writers have a blocking interface to read or write the next item, which
suffices for most simple applications. The vertex can report status and
errors to the job manager, and the progress of channels is automatically
monitored.

Many developers find it convenient to inherit from predefined vertex
classes that hide the details of the underlying channels and vertices.
We supply map and reduce classes with similar interfaces to those
described in {[}16{]}. We have also written a variety of others
including a general-purpose distribute that takes a single input stream
and writes on multiple outputs, and joins that call a virtual method
with every matching record tuple. These ``classes'' are simply vertices
like any other, so it is straightforward to write new ones to support
developers working in a particular domain.

\section{4.2 Legacy executables}\label{legacy-executables}

We provide a library ``process wrapper'' vertex that forks an executable
supplied as an invocation parameter. The wrapper vertex must work with
arbitrary data types, so its ``items'' are simply fixed-size buffers
that are passed unmodified to the forked process using named pipes in
the file-system. This allows unmodified pre-existing binaries to be run
as Dryad vertex programs. It is easy, for example, to invoke perl
scripts or grep at some vertices of a Dryad job.

\section{4.3 Efficient pipelined
execution}\label{efficient-pipelined-execution}

Most Dryad vertices contain purely sequential code. We also support an
event-based programming style, using a shared thread pool. The program
and channel interfaces have asynchronous forms, though unsurprisingly it
is harder to use the asynchronous interfaces than it is to write
sequential code using the synchronous interfaces. In some cases it may
be worth investing this effort, and many of the standard Dryad vertex
classes, including non-deterministic merge, sort, and generic maps and
joins, are built using the event-based programming style. The runtime
automatically distinguishes between vertices which can use a thread pool
and those that require a dedicated thread, and therefore encapsulated
graphs which contain hundreds of asynchronous vertices are executed
efficiently on a shared thread pool.

The channel implementation schedules read, write, serialization and
deserialization tasks on a thread pool shared between all channels in a
process, and a vertex can concurrently read or write on hundreds of
channels. The runtime tries to ensure efficient pipelined execution
while still presenting the developer with the simple abstraction of
reading and writing a single record at a time. Extensive use is made of
batching {[}28{]} to try to ensure that threads process hundreds or
thousands of records at a time without touching a reference count or
accessing a shared queue. The experiments in Section 6.2 substantiate
our claims for the efficiency of these abstractions: even single-node
Dryad applications have throughput comparable to that of a commercial
database system.

\section{5. JOB EXECUTION}\label{job-execution}

The scheduler inside the job manager keeps track of the state and
history of each vertex in the graph. At present if the job manager's
computer fails the job is terminated, though the vertex scheduler could
employ checkpointing or replication to avoid this. A vertex may be
executed multiple times over the length of the job due to failures, and
more than one instance of a given vertex may be executing at any given
time. Each execution of the vertex has a version number and a
corresponding ``execution record'' that contains the state of that
execution and the versions of the predecessor vertices from which its
inputs are derived. Each execution names its file-based output channels
uniquely using its version number to avoid conflicts among versions. If
the entire job completes successfully then each vertex selects a
successful execution and renames its output files to their correct final
forms.

When all of a vertex's input channels become ready a new execution
record is created for the vertex and placed in a scheduling queue. A
disk-based channel is considered to be ready when the entire file is
present. A channel that is a TCP pipe or shared-memory FIFO is ready
when the predecessor vertex has at least one running execution record.

A vertex and any of its channels may each specify a ``hard-constraint''
or a ``preference'' listing the set of computers on which it would like
to run. The constraints are combined and attached to the execution
record when it is added to the scheduling queue and they allow the
application writer to require that a vertex be co-located with a large
input file, and in general let the scheduler preferentially run
computations close to their data.

At present the job manager performs greedy scheduling based on the
assumption that it is the only job running on the cluster. When an
execution record is paired with an available computer the remote daemon
is instructed to run the specified vertex, and during execution the job
manager receives periodic status updates from the vertex. If every
vertex eventually completes then the job is deemed to have completed
successfully. If any vertex is re-run more than a set number of times
then the entire job is failed.

Files representing temporary channels are stored in directories managed
by the daemon and cleaned up after the job completes, and vertices are
killed by the daemon if their ``parent'' job manager crashes. We have a
simple graph visualizer suitable for small jobs that shows the state of
each vertex and the amount of data transmitted along each channel as the
computation progresses. A web-based interface shows regularly-updated
summary statistics of a running job and can be used to monitor large
computations. The statistics include the number of vertices that have
completed or been re-executed, the amount of data transferred across
channels, and the error codes reported by failed vertices. Links are
provided from the summary page that allow a developer to download logs
or crash dumps for further debugging, along with a script that allows
the vertex to be re-executed in isolation on a local machine.

\section{5.1 Fault tolerance policy}\label{fault-tolerance-policy}

Failures are to be expected during the execution of any distributed
application. Our default failure policy is suitable for the common case
that all vertex programs are deterministic. \(^{1}\) Because our
communication graph is acyclic, it is relatively straightforward to
ensure that every terminating execution of a job with immutable inputs
will compute the

same result, regardless of the sequence of computer or disk failures
over the course of the execution.

When a vertex execution fails for any reason the job manager is
informed. If the vertex reported an error cleanly the process forwards
it via the daemon before exiting; if the process crashes the daemon
notifies the job manager; and if the daemon fails for any reason the job
manager receives a heartbeat timeout. If the failure was due to a read
error on an input channel (which is be reported cleanly) the default
policy also marks the execution record that generated that version of
the channel as failed and terminates its process if it is running. This
will cause the vertex that created the failed input channel to be
re-executed, and will lead in the end to the offending channel being
re-created. Though a newly-failed execution record may have non-failed
successor records, errors need not be propagated forwards: since
vertices are deterministic two successors may safely compute using the
outputs of different execution versions. Note however that under this
policy an entire connected component of vertices connected by pipes or
shared-memory FIFOs will fail as a unit since killing a running vertex
will cause it to close its pipes, propagating errors in both directions
along those edges. Any vertex whose execution record is set to failed is
immediately considered for re-execution.

As Section 3.6 explains, each vertex belongs to a ``stage,'' and each
stage has a manager object that receives a callback on every state
transition of a vertex execution in that stage, and on a regular timer
interrupt. Within this callback the stage manager holds a global lock on
the job manager data-structures and can therefore implement quite
sophisticated behaviors. For example, the default stage manager includes
heuristics to detect vertices that are running slower than their peers
and schedule duplicate executions. This prevents a single slow computer
from delaying an entire job and is similar to the backup task mechanism
reported in {[}16{]}. In future we may allow non-deterministic vertices,
which would make fault-tolerance more interesting, and so we have
implemented our policy via an extensible mechanism that allows
non-standard applications to customize their behavior.

\section{5.2 Run-time graph refinement}\label{run-time-graph-refinement}

We have used the stage-manager callback mechanism to implement run-time
optimization policies that allow us to scale to very large input sets
while conserving scarce network bandwidth. Some of the large clusters we
have access to have their network provisioned in a two-level hierarchy,
with a dedicated mini-switch serving the computers in each rack, and the
per-rack switches connected via a single large core switch. Therefore
where possible it is valuable to schedule vertices as much as possible
to execute on the same computer or within the same rack as their input
data.

If a computation is associative and commutative, and performs a data
reduction, then it can benefit from an aggregation tree. As shown in
Figure 6, a logical graph connecting a set of inputs to a single
downstream vertex can be refined by inserting a new layer of internal
vertices, where each internal vertex reads data from a subset of the
inputs that are close in network topology, for example on the same
computer or within the same rack. If the internal vertices perform a
data reduction, the overall network traffic between racks will be
reduced by this refinement. A typical application would be a
histogramming operation that takes as input a set of partial histograms
and outputs their union. The implementation in Dryad simply attaches a
custom stage manager to the input layer. As this aggregation manager
receives callback notifications that upstream vertices have completed,
it rewrites the graph with the appropriate refinements.

Before\\
\textit{[Image omitted: images/92edfc0d3f8b8b7308dc95a06b229577ab42afd0080bd6996b26b241d60f1936.jpg]}

flowchart

\begin{Shaded}
\begin{Highlighting}[]
\NormalTok{graph TD}
\NormalTok{    B[" B "] {-}{-}\textgreater{} A1["A"]}
\NormalTok{    B {-}{-}\textgreater{} A2["A"]}
\NormalTok{    B {-}{-}\textgreater{} A3["A"]}
\NormalTok{    B {-}{-}\textgreater{} A4["A"]}
\NormalTok{    B {-}{-}\textgreater{} A5["A"]}
\end{Highlighting}
\end{Shaded}

After\\
\textit{[Image omitted: images/d9b79ed44567f8efbfa222b30fb9637a31bfb6042c1b3552ba7ea9876cdf45cb.jpg]}

flowchart

\begin{Shaded}
\begin{Highlighting}[]
\NormalTok{graph TD}
\NormalTok{    A["•"] {-}{-}\textgreater{} B["•"]}
\NormalTok{    A {-}{-}\textgreater{} Z["+"]}
\NormalTok{    B {-}{-}\textgreater{} Z}
\NormalTok{    Z {-}{-}\textgreater{} A}
\NormalTok{    Z {-}{-}\textgreater{} B}
\NormalTok{    Z {-}{-}\textgreater{} Z}
\NormalTok{    Z {-}{-}\textgreater{} A}
\NormalTok{    A {-}{-}\textgreater{} Z}
\NormalTok{    B {-}{-}\textgreater{} Z}
\NormalTok{    Z {-}{-}\textgreater{} A}
\NormalTok{    style Z stroke{-}dasharray: 5 5}
\end{Highlighting}
\end{Shaded}

Figure 6: A dynamic refinement for aggregation. The logical graph on the
left connects every input to the single output. The locations and sizes
of the inputs are not known until run time when it is determined which
computer each vertex is scheduled on. At this point the inputs are
grouped into subsets that are close in network topology, and an internal
vertex is inserted for each subset to do a local aggregation, thus
saving network bandwidth. The internal vertices are all of the same
user-supplied type, in this case shown as ``Z.'' In the diagram on the
right, vertices with the same label (`+' or `*') are executed close to
each other in network topology.

The operation in Figure 6 can be performed recursively to generate as
many layers of internal vertices as required. We have also found a
``partial aggregation'' operation to be very useful. This refinement is
shown in Figure 7; having grouped the inputs into k sets, the optimizer
replicates the downstream vertex k times to allow all of the sets to be
processed in parallel. Optionally, the partial refinement can be made to
propagate through the graph so that an entire pipeline of vertices will
be replicated k times (this behavior is not shown in the figure). An
example of the application of this technique is described in the
experiments in Section 6.3. Since the aggregation manager is notified on
the completion of upstream vertices, it has access to the size of the
data written by those vertices as well as its location. A typical
grouping heuristic ensures that a downstream vertex has no more than a
set number of input channels, or a set volume of input data. A special
case of partial refinement can be performed at startup to size the
initial layer of a graph so that, for example, each vertex processes
multiple inputs up to some threshold with the restriction that all the
inputs must lie on the same computer.

Because input data can be replicated on multiple computers in a cluster,
the computer on which a graph vertex is scheduled is in general
non-deterministic. Moreover the

Before\\
\textit{[Image omitted: images/19a2607b170426ce3048bcfdd65a7169d049adc0d96485f03f15ed4281c1dd48.jpg]}

flowchart

\begin{Shaded}
\begin{Highlighting}[]
\NormalTok{graph TD}
\NormalTok{    A["A"] {-}{-}\textgreater{} B["B"]}
\NormalTok{    C["C"] {-}{-}\textgreater{} B["B"]}
\NormalTok{    A["A"] {-}{-}\textgreater{} B["B"]}
\NormalTok{    A["A"] {-}{-}\textgreater{} B["B"]}
\NormalTok{    A["A"] {-}{-}\textgreater{} B["B"]}
\NormalTok{    A["A"] {-}{-}\textgreater{} B["B"]}
\NormalTok{    A["A"] {-}{-}\textgreater{} B["B"]}
\NormalTok{    B["B"] {-}{-}\textgreater{} C["C"]}
\NormalTok{    B["B"] {-}{-}\textgreater{} C["C"]}
\NormalTok{    B["B"] {-}{-}\textgreater{} A["A"]}
\NormalTok{    B["B"] {-}{-}\textgreater{} A["A"]}
\NormalTok{    B["B"] {-}{-}\textgreater{} A["A"]}
\NormalTok{    B["B"] {-}{-}\textgreater{} A["A"]}
\NormalTok{    B["B"] {-}{-}\textgreater{} A["A"]}
\end{Highlighting}
\end{Shaded}

After\\
\textit{[Image omitted: images/b159fad0c30036cd4697dbedd9549650dcdcdc8696b9ecb913fa68bcd4d6e31e.jpg]}

flowchart

Neural network architecture diagram showing connections between input
nodes A, B, C and output nodes with a central node marked k

Figure 7: A partial aggregation refinement. Following an input grouping
as in Figure 6 into k sets, the successor vertex is replicated k times
to process all the sets in parallel.

amount of data written in intermediate computation stages is typically
not known before a computation begins. Therefore dynamic refinement is
often more efficient than attempting a static grouping in advance.

Dynamic refinements of this sort emphasize the power of overlaying a
physical graph with its ``skeleton.'' For many applications, there is an
equivalence class of graphs with the same skeleton that compute the same
result. Varying the number of vertices in each stage, or their
connectivity, while preserving the graph topology at the stage level, is
merely a (dynamic) performance optimization.

\section{6. EXPERIMENTAL EVALUATION}\label{experimental-evaluation}

Dryad has been used for a wide variety of applications, including
relational queries, large-scale matrix computations, and many
text-processing tasks. For this paper we examined the effectiveness of
the Dryad system in detail by running two sets of experiments. The first
experiment takes the SQL query described in Section 2.1 and implements
it as a Dryad application. We compare the Dryad performance with that of
a traditional commercial SQL server, and we analyze the Dryad
performance as the job is distributed across different numbers of
computers. The second is a simple map-reduce style data-mining
operation, expressed as a Dryad program and applied to 10.2 TBytes of
data using a cluster of around 1800 computers.

The strategies we adopt to build our communication flow graphs are
familiar from the parallel database literature \([18]\) and include
horizontally partitioning the datasets, exploiting pipelined parallelism
within processes and applying exchange operations to communicate partial
results between the partitions. None of the application-level code in
any of our experiments makes explicit use of concurrency primitives.

\section{6.1 Hardware}\label{hardware}

The SQL query experiments were run on a cluster of 10 computers in our
own laboratory, and the data-mining tests were run on a cluster of
around 1800 computers embedded in a data center. Our laboratory
computers each had 2 dual-core Opteron processors running at 2 GHz
(i.e., 4 CPUs total), 8 GBytes of DRAM (half attached to each processor
chip), and 4 disks. The disks were 400 GByte Western Digital WD40
00YR-01PLB0 SATA drives, connected through a Silicon Image 3114 PCI SATA
controller (66MHz, 32-bit). Network connectivity was by 1 Gbit/sec
Ethernet links connecting into a single non-blocking switch. One of our
laboratory computers was dedicated to running SQLServer and its data was
stored in 4 separate 350 GByte NTFS volumes, one on each drive, with
SQLServer configured to do its own data striping for the raw data and
for its temporary tables. All the other laboratory computers were
configured with a single 1.4 TByte NTFS volume on each computer, created
by software striping across the 4 drives. The computers in the data
center had a variety of configurations, but were typically roughly
comparable to our laboratory equipment. All the computers were running
Windows Server 2003 Enterprise x64 edition SP1.

\section{6.2 SQL Query}\label{sql-query}

The query for this experiment is described in Section 2.1 and uses the
Dryad communication graph shown in Figure 2. SQLServer 2005's execution
plan for this query was very close to the Dryad computation, except that
it used an external hash join for ``Y'' in place of the sort-merge we
chose for Dryad. SQLServer takes slightly longer if it is forced by a
query hint to use a sort-merge join.

For our experiments, we used two variants of the Dryad graph:
``in-memory'' and ``two-pass.'' In both variants communication from each
\(M_{i}\) through its corresponding \(S_{i}\) to Y is by a shared-memory
FIFO. This pulls four sorters into the same process to execute in
parallel on the four CPUs in each computer. In the ``in-memory'' variant
only, communication from each \(D_{i}\) to its four corresponding
\(M_{j}\) vertices is also by a shared-memory FIFO and the rest of the
\(D_{i} \rightarrow M_{k}\) edges use TCP pipes. All other communication
is through NTFS temporary files in both variants.

There is good spatial locality in the query, which improves as the
number of partitions (n) decreases: for n = 40 an average of 80\% of the
output of \(D_{i}\) goes to its corresponding \(M_{i}\) , increasing to
88\% for n = 6. In either variant n must be large enough that every sort
executed by a vertex \(S_{i}\) will fit into the computer's 8 GBytes of
DRAM (or else it will page). With the current data, this threshold is at
n = 6.

Note that the non-deterministic merge in M randomly permutes its output
depending on the order of arrival of items on its input channels and
this technically violates the requirement that all vertices be
deterministic. This does not cause problems for our fault-tolerance
model because the sort \(S_{i}\) ``undoes'' this permutation, and since
the edge from \(M_{i}\) to \(S_{i}\) is a shared-memory FIFO within a
single process the two vertices fail (if at all) in tandem and the
non-determinism never ``escapes.''

The in-memory variant requires at least n computers since otherwise the
S vertices will deadlock waiting for data from an X vertex. The two-pass
variant will run on any number of computers. One way to view this
trade-off is that by adding the file buffering in the two-pass variant
we in effect converted to using a two-pass external sort. Note that the
conversion from the in-memory to the two-pass program simply involves
changing two lines in the graph construction code, with no modifications
to the vertex programs.

We ran the two-pass variant using n = 40, varying the number of
computers from 1 to 9. We ran the in-memory variant using n = 6 through
n = 9, each time on n computers. As a baseline measurement we ran the
query on a reasonably well optimized SQLServer on one computer. Table 2
shows the elapsed times in seconds for each experiment. On repeated runs
the times were consistent to within 3.4\% of their averages except for
the single-computer two-pass case, which was within 9.4\%. Figure 8
graphs the inverse of these times, normalized to show the speed-up
factor relative to the two-pass single-computer case.

The results are pleasantly straightforward. The two-pass Dryad job works
on all cluster sizes, with close to linear speed-up. The in-memory
variant works as expected for

Computers

1

2

3

4

5

6

7

8

9

SQLServer

3780

Two-pass

2370

1260

836

662

523

463

423

346

321

In-memory

217

203

183

168

Table 2: Time in seconds to process an SQL query using different numbers
of computers. The SQLServer implementation cannot be distributed across
multiple computers and the in-memory experiment can only be run for 6 or
more computers.

\textit{[Image omitted: images/81af627398c0df4669b6e39a34e16bfd0cabc6861994c1391503872baa325e74.jpg]}

line

{\def\LTcaptype{none} % do not increment counter
\begin{longtable}[]{@{}
  >{\raggedright\arraybackslash}p{(\linewidth - 6\tabcolsep) * \real{0.3065}}
  >{\raggedright\arraybackslash}p{(\linewidth - 6\tabcolsep) * \real{0.2419}}
  >{\raggedright\arraybackslash}p{(\linewidth - 6\tabcolsep) * \real{0.2258}}
  >{\raggedright\arraybackslash}p{(\linewidth - 6\tabcolsep) * \real{0.2258}}@{}}
\toprule\noalign{}
\begin{minipage}[b]{\linewidth}\raggedright
Number of Computers
\end{minipage} & \begin{minipage}[b]{\linewidth}\raggedright
Dryad In-Memory
\end{minipage} & \begin{minipage}[b]{\linewidth}\raggedright
Dryad Two-Pass
\end{minipage} & \begin{minipage}[b]{\linewidth}\raggedright
SQLServer 2005
\end{minipage} \\
\midrule\noalign{}
\endhead
\bottomrule\noalign{}
\endlastfoot
1 & - & 1.0 & 0.5 \\
2 & - & 2.0 & - \\
4 & - & 3.5 & - \\
6 & 11.0 & 5.0 & - \\
8 & 13.0 & 7.0 & - \\
9 & 14.0 & 7.5 & - \\
\end{longtable}
}

Figure 8: The speedup of the SQL query computation is near-linear in the
number of computers used. The baseline is relative to Dryad running on a
single computer and times are given in Table 2.

n = 6 and up, again with close to linear speed-up, and approximately
twice as fast as the two-pass variant. The SQLServer result matches our
expectations: our specialized Dryad program runs significantly, but not
outrageously, faster than SQLServer's general-purpose query engine. We
should note of course that Dryad simply provides an execution engine
while the database provides much more functionality, including logging,
transactions, and mutable relations.

\section{6.3 Data mining}\label{data-mining}

The data-mining experiment fits the pattern of map then reduce. The
purpose of running this experiment was to verify that Dryad works
sufficiently well in these straightforward cases, and that it works at
large scales.

The computation in this experiment reads query logs gathered by the MSN
Search service, extracts the query strings, and builds a histogram of
query frequency. The basic communication graph is shown in Figure 9. The
log files are partitioned and replicated across the computers' disks.
The P vertices each read their part of the log files using library
newline-delimited text items, and parse them to extract the query
strings. Subsequent items are all library tuples containing a query
string, a count, and a hash of the string. Each D vertex distributes to
k outputs based on the query string hash; S performs an in-memory sort.
C accumulates total counts for each query and MS performs a streaming
merge-sort. S and MS come from a vertex library and take a comparison
function as a parameter; in this example they sort based on the query
hash. We have encapsulated the simple vertices into subgraphs denoted by
diamonds in order to reduce the total number of vertices in the job (and
hence the overhead associated with process start-up) and the volume of
temporary data written to disk.

The graph shown in Figure 9 does not scale well to very large datasets.
It is wasteful to execute a separate Q vertex for every input partition.
Each partition is only around 100 MBytes, and the P vertex performs a
substantial data reduction, so the amount of data which needs to be
sorted by the S vertices is very much less than the total RAM on a
computer. Also, each R subgraph has n inputs, and when n grows to
hundreds of thousands of partitions, it becomes unwieldy to read in
parallel from so many channels.

\textit{[Image omitted: images/6a5462d1def762a4f4aa325aabe87a0aaaa147f8d9e607b67084d45edb1e3b69.jpg]}

flowchart

\begin{Shaded}
\begin{Highlighting}[]
\NormalTok{graph TD}
\NormalTok{    subgraph LeftPanel}
\NormalTok{        R1["Square"] {-}{-}\textgreater{} R2["Diamond"]}
\NormalTok{        R2 {-}{-}\textgreater{} R3["Square"]}
\NormalTok{        R3 {-}{-}\textgreater{} Q1["Q"]}
\NormalTok{        R1 {-}{-}\textgreater{} R4["Diamond"]}
\NormalTok{        R4 {-}{-}\textgreater{} R5["Square"]}
\NormalTok{        R5 {-}{-}\textgreater{} Q4["Q"]}
\NormalTok{        R2 {-}{-}\textgreater{} R6["Square"]}
\NormalTok{        R6 {-}{-}\textgreater{} R7["Square"]}
\NormalTok{        R1 {-}{-}\textgreater{} R8["Square"]}
\NormalTok{        R8 {-}{-}\textgreater{} R9["Square"]}
\NormalTok{        R2 {-}{-}\textgreater{} R10["Square"]}
\NormalTok{        R10 {-}{-}\textgreater{} R11["Square"]}
\NormalTok{        R3 {-}{-}\textgreater{} R12["Square"]}
\NormalTok{        R12 {-}{-}\textgreater{} R13["Square"]}
\NormalTok{        R4 {-}{-}\textgreater{} R14["Square"]}
\NormalTok{        R14 {-}{-}\textgreater{} R15["Square"]}
\NormalTok{        R5 {-}{-}\textgreater{} R16["Square"]}
\NormalTok{        R16 {-}{-}\textgreater{} R17["Square"]}
\NormalTok{        R2 {-}{-}\textgreater{} R18["Square"]}
\NormalTok{        R18 {-}{-}\textgreater{} R19["Square"]}
\NormalTok{        R3 {-}{-}\textgreater{} R20["Square"]}
\NormalTok{        R20 {-}{-}\textgreater{} R21["Square"]}
\NormalTok{        R4 {-}{-}\textgreater{} R22["Square"]}
\NormalTok{        R22 {-}{-}\textgreater{} R23["Square"]}
\NormalTok{        R5 {-}{-}\textgreater{} R24["Square"]}
\NormalTok{        R24 {-}{-}\textgreater{} R25["Square"]}
\NormalTok{        R6 {-}{-}\textgreater{} R26["Square"]}
\NormalTok{        R26 {-}{-}\textgreater{} R27["Square"]}
\NormalTok{        R7 {-}{-}\textgreater{} R28["Square"]}
\NormalTok{        R28 {-}{-}\textgreater{} R29["Square"]}
\NormalTok{        R3 {-}{-}\textgreater{} R30["Square"]}
\NormalTok{        R30 {-}{-}\textgreater{} R31["Square"]}
\NormalTok{        R4 {-}{-}\textgreater{} R32["Square"]}
\NormalTok{        R32 {-}{-}\textgreater{} R33["Square"]}
\NormalTok{        R5 {-}{-}\textgreater{} R34["Square"]}
\NormalTok{        R34 {-}{-}\textgreater{} R35["Square"]}
\NormalTok{        R6 {-}{-}\textgreater{} R36["Square"]}
\NormalTok{        R36 {-}{-}\textgreater{} R37["Square"]}
\NormalTok{        R7 {-}{-}\textgreater{} R38["Square"]}
\NormalTok{        R38 {-}{-}\textgreater{} R39["Square"]}
\NormalTok{        R8 {-}{-}\textgreater{} R40["Square"]}
\NormalTok{        R30 {-}{-}\textgreater{} R41["Square"]}
\NormalTok{        R31 {-}{-}\textgreater{} R42["Square"]}
\NormalTok{        R4 {-}{-}\textgreater{} R43["Square"]}
\NormalTok{        R5 {-}{-}\textgreater{} R44["Square"]}
\NormalTok{        R32 {-}{-}\textgreater{} R45["Square"]}
\NormalTok{        R33 {-}{-}\textgreater{} R46["Square"]}
\NormalTok{        R4 {-}{-}\textgreater{} R47["Square"]}
\NormalTok{        R5 {-}{-}\textgreater{} R48["Square"]}
\NormalTok{        R33 {-}{-}\textgreater{} R49["Square"]}
\NormalTok{        R34 {-}{-}\textgreater{} R50["Square"]}
\NormalTok{        R4 {-}{-}\textgreater{} R51["Square"]}
\NormalTok{        R5 {-}{-}\textgreater{} R52["Square"]}
\NormalTok{        R6 {-}{-}\textgreater{} R53["Square"]}
\NormalTok{        R33 {-}{-}\textgreater{} R54["Square"]}
\NormalTok{        R5 {-}{-}\textgreater{} R55["Square"]}
\NormalTok{        R6 {-}{-}\textgreater{} R56["Square"]}
\NormalTok{        R34 {-}{-}\textgreater{} R57["Square"]}
\NormalTok{        R5 {-}{-}\textgreater{} R58["Square"]}
\NormalTok{        R6 {-}{-}\textgreater{} R59["Square"]}
\NormalTok{        R35 {-}{-}\textgreater{} R60["Square"]}
\NormalTok{        R5 {-}{-}\textgreater{} R61["Square"]}
\NormalTok{        R6 {-}{-}\textgreater{} R62["Square"]}
\NormalTok{        R7 {-}{-}\textgreater{} R63["Square"]}
\NormalTok{        R36 {-}{-}\textgreater{} R64["Square"]}
\NormalTok{        R5 {-}{-}\textgreater{} R65["Square"]}
\NormalTok{        R6 {-}{-}\textgreater{} R66["Square"]}
\NormalTok{        R7 {-}{-}\textgreater{} R67["Square"]}
\NormalTok{        R37 {-}{-}\textgreater{} R68["Square"]}
\NormalTok{        R5 {-}{-}\textgreater{} R69["Square"]}
\NormalTok{        R6 {-}{-}\textgreater{} R70["Square"]}
\NormalTok{        R7 {-}{-}\textgreater{} R71["Square"]}
\NormalTok{        R8 {-}{-}\textgreater{} R72["Square"]}
\NormalTok{        R38 {-}{-}\textgreater{} R73["Square"]}
\NormalTok{        R5 {-}{-}\textgreater{} R74["Square"]}
\NormalTok{        R6 {-}{-}\textgreater{} R75["Square"]}
\NormalTok{        R7 {-}{-}\textgreater{} R76["Square"]}
\NormalTok{        R8 {-}{-}\textgreater{} R77["Square"]}
\NormalTok{        R38 {-}{-}\textgreater{} R78["Square"]}
\NormalTok{        R5 {-}{-}\textgreater{} R79["Square"]}
\NormalTok{        R6 {-}{-}\textgreater{} R80["Square"]}
\NormalTok{        R7 {-}{-}\textgreater{} R81["Square"]}
\NormalTok{        R8 {-}{-}\textgreater{} R82["Square"]}
\NormalTok{        R9 {-}{-}\textgreater{} R83["Square"]}
\NormalTok{        R39 {-}{-}\textgreater{} R84["Square"]}
\NormalTok{        R5 {-}{-}\textgreater{} R85["Square"]}
\NormalTok{        R6 {-}{-}\textgreater{} R86["Square"]}
\NormalTok{        R7 {-}{-}\textgreater{} R87["Square"]}
\NormalTok{        R8 {-}{-}\textgreater{} R88["Square"]}
\NormalTok{        R9 {-}{-}\textgreater{} R89["Square"]}
\NormalTok{        R10 {-}{-}\textgreater{} R90["Square"]}
\NormalTok{        R11 {-}{-}\textgreater{} R91["Square"]}
\NormalTok{        R12 {-}{-}\textgreater{} R92["Square"]}
\NormalTok{        R13 {-}{-}\textgreater{} R93["Square"]}
\NormalTok{        R14 {-}{-}\textgreater{} R94["Square"]}
\NormalTok{        R15 {-}{-}\textgreater{} R95["Square"]}
\NormalTok{        R16 {-}{-}\textgreater{} R96["Square"]}
\NormalTok{        R17 {-}{-}\textgreater{} R97["Square"]}
\NormalTok{        R18 {-}{-}\textgreater{} R98["Square"]}
\NormalTok{        R19 {-}{-}\textgreater{} R99["Square"]}
\NormalTok{        R20 {-}{-}\textgreater{} R100["Square"]}
\NormalTok{        R21 {-}{-}\textgreater{} R101["Square"]}
\NormalTok{        R22 {-}{-}\textgreater{} R102["Square"]}
\NormalTok{        R23 {-}{-}\textgreater{} R103["Square"]}
\NormalTok{        R24 {-}{-}\textgreater{} R104["Square"]}
\NormalTok{        R25 {-}{-}\textgreater{} R105["Square"]}
\NormalTok{        R26 {-}{-}\textgreater{} R106["Square"]}
\NormalTok{        R27 {-}{-}\textgreater{} R107["Square"]}
\NormalTok{        R28 {-}{-}\textgreater{} R108["Square"]}
\NormalTok{        R29 {-}{-}\textgreater{} R109["Square"]}
\NormalTok{        R30 {-}{-}\textgreater{} R110["Square"]}
\NormalTok{        R31 {-}{-}\textgreater{} R111["Square"]}
\NormalTok{        R32 {-}{-}\textgreater{} R112["Square"]}
\NormalTok{        R33 {-}{-}\textgreater{} R113["Square"]}
\NormalTok{        R34 {-}{-}\textgreater{} R114["Square"]}
\NormalTok{        R35 {-}{-}\textgreater{} R115["Square"]}
\NormalTok{        R36 {-}{-}\textgreater{} R116["Square"]}
\NormalTok{        R37 {-}{-}\textgreater{} R117["Square"]}
\NormalTok{        R38 {-}{-}\textgreater{} R118["Square"]}
\NormalTok{        R39 {-}{-}\textgreater{} R119["Square"]}
\NormalTok{        R40 {-}{-}\textgreater{} R120["Square"]}
\NormalTok{        R41 {-}{-}\textgreater{} R121["Square"]}
\NormalTok{        R42 {-}{-}\textgreater{} R122["Square"]}
\NormalTok{        R43 {-}{-}\textgreater{} R123["Square"]}
\NormalTok{        R44 {-}{-}\textgreater{} R124["Square"]}
\NormalTok{        R45 {-}{-}\textgreater{} R125["Square"]}
\NormalTok{        R46 {-}{-}\textgreater{} R126["Square"]}
\NormalTok{        R47 {-}{-}\textgreater{} R127["Square"]}
\NormalTok{        R48 {-}{-}\textgreater{} R128["Square"]}
\NormalTok{        R49 {-}{-}\textgreater{} R129["Square"]}
\NormalTok{        R50 {-}{-}\textgreater{} R130["Square"]}
\NormalTok{        R51 {-}{-}\textgreater{} R131["Square"]}
\NormalTok{        R52 {-}{-}\textgreater{} R132["Square"]}
\NormalTok{        R53 {-}{-}\textgreater{} R133["Square"]}
\NormalTok{        R54 {-}{-}\textgreater{} R134["Square"]}
\NormalTok{        R55 {-}{-}\textgreater{} R135["Square"]}
\NormalTok{        R56 {-}{-}\textgreater{} R136["Square"]}
\NormalTok{        R57 {-}{-}\textgreater{} R137["Square"]}
\NormalTok{        R58 {-}{-}\textgreater{} R138["Square"]}
\NormalTok{        R59 {-}{-}\textgreater{} R139["Square"]}
\NormalTok{        R60 {-}{-}\textgreater{} R140["Square"]}
\NormalTok{        R61 {-}{-}\textgreater{} R141["Square"]}
\NormalTok{        R62 {-}{-}\textgreater{} R142["Square"]}
\NormalTok{        R63 {-}{-}\textgreater{} R143["Square"]}
\NormalTok{        R64 {-}{-}\textgreater{} R144["Square"]}
\NormalTok{        R65 {-}{-}\textgreater{} R145["Square"]}
\NormalTok{        R66 {-}{-}\textgreater{} R146["Square"]}
\NormalTok{        R67 {-}{-}\textgreater{} R147["Square"]}
\NormalTok{        R68 {-}{-}\textgreater{} R148["Square"]}
\NormalTok{        R69 {-}{-}\textgreater{} R149["Square"]}
\NormalTok{        R70 {-}{-}\textgreater{} R150["Square"]}
\NormalTok{        R71 {-}{-}\textgreater{} R151["Square"]}
\NormalTok{        R72 {-}{-}\textgreater{} R152["Square"]}
\NormalTok{        R73 {-}{-}\textgreater{} R153["Square"]}
\NormalTok{        R74 {-}{-}\textgreater{} R154["Square"]}
\NormalTok{        R75 {-}{-}\textgreater{} R155["Square"]}
\NormalTok{        R76 {-}{-}\textgreater{} R156["Square"]}
\NormalTok{        R77 {-}{-}\textgreater{} R157["Square"]}
\NormalTok{        R78 {-}{-}\textgreater{} R158["Square"]}
\NormalTok{        R79 {-}{-}\textgreater{} R159["Square"]}
\NormalTok{        R80 {-}{-}\textgreater{} R160["Square"]}
\NormalTok{        R81 {-}{-}\textgreater{} R161["Square"]}
\NormalTok{        R82 {-}{-}\textgreater{} R162["Square"]}
\NormalTok{        R83 {-}{-}\textgreater{} R163["Square"]}
\NormalTok{        R84 {-}{-}\textgreater{} R164["Square"]}
\NormalTok{        R85 {-}{-}\textgreater{} R165["Square"]}
\NormalTok{        R86 {-}{-}\textgreater{} R166["Square"]}
\NormalTok{        R87 {-}{-}\textgreater{} R167["Square"]}
\NormalTok{        R88 {-}{-}\textgreater{} R168["Square"]}
\NormalTok{        R89 {-}{-}\textgreater{} R169["Square"]}
\NormalTok{        R90 {-}{-}\textgreater{} R170["Square"]}
\NormalTok{        R91 {-}{-}\textgreater{} R171["Square"]}
\NormalTok{        R92 {-}{-}\textgreater{} R172["Square"]}
\NormalTok{        R93 {-}{-}\textgreater{} R173["Square"]}
\NormalTok{        R94 {-}{-}\textgreater{} R174["Square"]}
\NormalTok{        R95 {-}{-}\textgreater{} R175["Square"]}
\NormalTok{        R96 {-}{-}\textgreater{} R176["Square"]}
\NormalTok{        R97 {-}{-}\textgreater{} R177["Square"]}
\NormalTok{        R98 {-}{-}\textgreater{} R178["Square"]}
\NormalTok{        R99 {-}{-}\textgreater{} R179["Square"]}
\NormalTok{        R100 {-}{-}\textgreater{} R180["Square"]}
\NormalTok{    end}
\NormalTok{    subgraph RightPanel}
\NormalTok{        Each["Each Q is:"]}
\NormalTok{        Each["Each R is:"]}
\NormalTok{        Each["Each C is:"]}
\NormalTok{        Each["Each R is:"]}
\NormalTok{        Each["Each C is:"]}
\NormalTok{        Each["Each MS is:"]}
\NormalTok{    end}
\end{Highlighting}
\end{Shaded}

Figure 9: The communication graph to compute a query histogram. Details
are in Section 6.3. This figure shows the first cut ``naive''
encapsulated version that doesn't scale well.

After trying a number of different encapsulation and dynamic refinement
schemes we arrived at the communication graphs shown in Figure 10 for
our experiment. Each subgraph in the first phase now has multiple
inputs, grouped automatically using the refinement in Figure 7 to ensure
they all lie on the same computer. The inputs are sent to the parser P
through a non-deterministic merge vertex M. The distribution (vertex D)
has been taken out of the first phase to allow another layer of grouping
and aggregation (again using the refinement in Figure 7) before the
explosion in the number of output channels.

We ran this experiment on 10,160,519,065,748 Bytes of input data in a
cluster of around 1800 computers embedded in a data center. The input
was divided into 99,713 partitions replicated across the computers, and
we specified that the application should use 450 R subgraphs. The first
phase grouped the inputs into at most 1 GBytes at a time, all lying on
the same computer, resulting in 10,405 \(Q'\) subgraphs that wrote a
total of 153,703,445,725 Bytes. The outputs from the \(Q'\) subgraphs
were then grouped into sets of at most 600 MBytes on the same local
switch resulting in 217 T subgraphs. Each T was connected to every R
subgraph, and they wrote a total of 118,364,131,628 Bytes. The total
output from the R subgraphs was 33,375,616,713 Bytes, and the end-to-end
computation took 11 minutes and 30 seconds. Though this experiment only
uses 11,072 vertices, intermediate experiments with other graph
topologies confirmed that Dryad can successfully execute jobs containing
hundreds of thousands of vertices.

We would like to emphasize several points about the optimization process
we used to arrive at the graphs in Figure 10:

\begin{enumerate}
\def\labelenumi{\arabic{enumi}.}
\tightlist
\item
  At no point during the optimization did we have to modify any of the
  code running inside the vertices: we were simply manipulating the
  graph of the job's communication flow, changing tens of lines of
  code.\\
\item
  This communication graph is well suited to any map-reduce computation
  with similar characteristics: i.e.~that the map phase (our \(P\)
  vertex) performs substan-
\end{enumerate}

\textit{[Image omitted: images/69cc853aed8934ce20208082b80442dc589641a621958b89f54d7edd04f27f02.jpg]}

flowchart

Diagram comparing two configurations (a and b) of a system with nodes
labeled R, T, C, and MS, showing data flow and memory usage patterns.

Figure 10: Rearranging the vertices gives better scaling performance
compared with Figure 9. The user supplies graph (a) specifying that 450
buckets should be used when distributing the output, and that each
\(Q'\) vertex may receive up to 1GB of input while each T may receive up
to 600MB. The number of \(Q'\) and T vertices is determined at run time
based on the number of partitions in the input and the network locations
and output sizes of preceding vertices in the graph, and the refined
graph (b) is executed by the system. Details are in Section 6.3.

tial data reduction and the reduce phase (our C vertex) performs some
additional relatively minor data reduction. A different topology might
give better performance for a map-reduce task with different behavior;
for example if the reduce phase performed substantial data reduction a
dynamic merge tree as described in Figure 6 might be more suitable.

\begin{enumerate}
\def\labelenumi{\arabic{enumi}.}
\setcounter{enumi}{2}
\tightlist
\item
  When scaling up another order of magnitude or two, we might change the
  topology again, e.g.~by adding more layers of aggregation between the
  T and R stages. Such re-factoring is easy to do.\\
\item
  Getting good performance for large-scale data-mining computations is
  not trivial. Many novel features of the Dryad system, including
  subgraph encapsulation and dynamic refinement, were used. These made
  it simple to experiment with different optimization schemes that would
  have been difficult or impossible to implement using a simpler but
  less powerful system.
\end{enumerate}

\section{7. BUILDING ON DRYAD}\label{building-on-dryad}

As explained in the introduction, we have targeted Dryad at developers
who are experienced at using high-level compiled programming languages.
In some domains there may be great value in making common large-scale
data processing tasks easier to perform, since this allows
non-developers to directly query the data store {[}33{]}. We designed
Dryad to be usable as a platform on which to develop such more
restricted but simpler programming interfaces, and two other groups
within Microsoft have already prototyped systems to address particular
application domains.

\section{7.1 The ``Nebula'' scripting
language}\label{the-nebula-scripting-language}

One team has layered a scripting interface on top of Dryad. It allows a
user to specify a computation as a series of stages (corresponding to
the Dryad stages described in Section 3.6), each taking inputs from one
or more previous stages or the file system. Nebula transforms Dryad into
a generalization of the Unix piping mechanism and it allows programmers
to write giant acyclic graphs spanning many computers. Often a Nebula
script only refers to existing executables such as perl or grep,
allowing a user to write an entire complex distributed application
without compiling any code. The Nebula layer on top of Dryad, together
with some perl wrapper functions, has proved to be very successful for
large-scale text processing, with a low barrier to entry for users.
Scripts typically run on thousands of computers and contain 5--15 stages
including multiple projections, aggregations and joins, often combining
the information from multiple input sets in sophisticated ways.

Nebula hides most of the details of the Dryad program from the
developer. Stages are connected to preceding stages using operators that
implicitly determine the number of vertices required. For example, a
``Filter'' operation creates one new vertex for every vertex in its
input list, and connects them pointwise to form a pipeline. An
``Aggregate'' operation can be used to perform exchanges and merges. The
implementation of the Nebula operators makes use of dynamic
optimizations like those described in Section 5.2 however the operator
abstraction allows users to remain unaware of the details of these
optimizations. All Nebula vertices execute the process wrapper described
in Section 4.2, and the vertices in a given stage all run the same
executable and command-line, specified using the script. The Nebula
system defines conventions for passing the names of the input and output
pipes to the vertex executable command-line.

There is a very popular ``front-end'' to Nebula that lets the user
describe a job using a combination of: fragments of perl that parse
lines of text from different sources into structured records; and a
relational query over those structured records expressed in a subset of
SQL that includes select, project and join. This job description is
converted into a Nebula script and executed using Dryad. The perl pars-

ing fragments for common input sources are all in libraries, so many
jobs using this front-end are completely described using a few lines of
SQL.

\section{7.2 Integration with SSIS}\label{integration-with-ssis}

SQL Server Integration Services (SSIS) {[}6{]} supports workflow-based
application programming on a single instance of SQLServer. The AdCenter
team in MSN has developed a system that embeds local SSIS computations
in a larger, distributed graph with communication, scheduling and fault
tolerance provided by Dryad. The SSIS input graph can be built and
tested on a single computer using the full range of SQL developer tools.
These include a graphical editor for constructing the job topology, and
an integrated debugger.

When the graph is ready to run on a larger cluster the system
automatically partitions it using heuristics and builds a Dryad graph
that is then executed in a distributed fashion. Each Dryad vertex is an
instance of SQLServer running an SSIS subgraph of the complete job. This
system is currently deployed in a live production system as part of one
of AdCenter's log processing pipelines.

\section{7.3 Distributed SQL queries}\label{distributed-sql-queries}

One obvious additional direction would be to adapt a query optimizer for
SQL or LINQ {[}4{]} queries to compile plans directly into a Dryad flow
graph using appropriate parameterized vertices for the relational
operations. Since our fault-tolerance model only requires that inputs be
immutable over the duration of the query, any underlying storage system
that offers lightweight snapshots would suffice to allow us to deliver
consistent query results. We intend to pursue this as future work.

\section{8. RELATED WORK}\label{related-work}

Dryad is related to a broad class of prior literature, ranging from
custom hardware to parallel databases, but we believe that the ensemble
of trade-offs we have chosen for its design, and some of the
technologies we have deployed, make it a unique system.

Hardware Several hardware systems use stream programming models similar
to Dryad, including Intel IXP {[}2{]}, Imagine {[}26{]}, and SCORE
{[}15{]}. Programmers or compilers represent the distributed computation
as a collection of independent subroutines residing within a high-level
graph.

Click A similar approach is adopted by the Click modular router
{[}27{]}. The technique used to encapsulate multiple Dryad vertices in a
single large vertex, described in section 3.4, is similar to the method
used by Click to group the elements (equivalent of Dryad vertices) in a
single process. However, Click is always single-threaded, while Dryad
encapsulated vertices are designed to take advantage of multiple CPU
cores that may be available.

Dataflow The overall structure of a Dryad application is closely related
to large-grain dataflow techniques used in e.g.~LGDF2 {[}19{]}, CODE2
{[}31{]} and P-RIO {[}29{]}. These systems were not designed to scale to
large clusters of commodity computers, however, and do not tolerate
machine failures or easily support programming very large graphs.
Paralex {[}9{]} has many similarities

to Dryad, but in order to provide automatic fault-tolerance sacrifices
the vertex programming model, allowing only pure-functional programs.

Parallel databases Dryad is heavily indebted to the traditional parallel
database field {[}18{]}: e.g., Vulcan {[}22{]}, Gamma {[}17{]}, RDb
{[}11{]}, DB2 parallel edition {[}12{]}, and many others. Many
techniques for exploiting parallelism, including data partitioning;
pipelined and partitioned parallelism; and hash-based distribution are
directly derived from this work. We can map the whole relational algebra
on top of Dryad, however Dryad is not a database engine: it does not
include a query planner or optimizer; the system has no concept of data
schemas or indices; and Dryad does not support transactions or logs.
Dryad gives the programmer more control than SQL via C++ programs in
vertices and allows programmers to specify encapsulation, transport
mechanisms for edges, and callbacks for vertex stages. Moreover, the
graph builder language allows Dryad to express irregular computations.

Continuous Query systems There are some superficial similarities between
CQ systems (e.g.~{[}25, 10, 34{]}) and Dryad, such as some operators and
the topologies of the computation networks. However, Dryad is a batch
computation system, not designed to support real-time operation which is
crucial for CQ systems since many CQ window operators depend on
real-time behavior. Moreover, many datamining Dryad computations require
extremely high throughput (tens of millions of records per second per
node), which is much greater than that typically seen in the CQ
literature.

Explicitly parallel languages like Parallel Haskell {[}38{]}, Cilk
{[}14{]} or NESL {[}13{]} have the same emphasis as Dryad on using the
user's knowledge of the problem to drive the parallelization. By relying
on C++, Dryad should have a faster learning curve than that for
functional languages, while also being able to leverage commercial
optimizing compilers. There is some appeal in these alternative
approaches, which present the user with a uniform programming
abstraction rather than our two-level hierarchy. However, we believe
that for data-parallel applications that are naturally written using
coarse-grain communication patterns, we gain substantial benefit by
letting the programmer cooperate with the system to decide on the
granularity of distribution.

Grid computing {[}1{]} and projects such as Condor {[}37{]} are clearly
related to Dryad, in that they leverage the resources of many
workstations using batch processing. However, Dryad does not attempt to
provide support for wide-area operation, transparent remote I/O, or
multiple administrative domains. Dryad is optimized for the case of a
very high-throughput LAN, whereas in Condor bandwidth management is
essentially handled by the user job.

Google MapReduce The Dryad system was primarily designed to support
large-scale data-mining over clusters of thousands of computers. As a
result, of the recent related systems it shares the most similarities
with Google's MapReduce {[}16, 33{]} which addresses a similar

problem domain. The fundamental difference between the two systems is
that a Dryad application may specify an arbitrary communication DAG
rather than requiring a sequence of map/distribute/sort/reduce
operations. In particular, graph vertices may consume multiple inputs,
and generate multiple outputs, of different types. For many applications
this simplifies the mapping from algorithm to implementation, lets us
build on a greater library of basic subroutines, and, together with the
ability to exploit TCP pipes and shared-memory for data edges, can bring
substantial performance gains. At the same time, our implementation is
general enough to support all the features described in the MapReduce
paper.

Scientific computing Dryad is also related to high-performance computing
platforms like MPI {[}5{]}, PVM {[}35{]}, or computing on GPUs {[}36{]}.
However, Dryad focuses on a model with no shared-memory between
vertices, and uses no synchronization primitives.

NOW The original impetus for employing clusters of workstations with a
shared-nothing memory model came from projects like Berkeley NOW {[}7,
8{]}, or TACC {[}20{]}. Dryad borrows some ideas from these systems,
such as fault-tolerance through re-execution and centralized resource
scheduling, but our system additionally provides a unified, simple
high-level programming language layer.

Log datamining Addamark, now renamed SenSage {[}32{]} has successfully
commercialized software for log datamining on clusters of workstations.
Dryad is designed to scale to much larger implementations, up to
thousands of computers.

\section{9. DISCUSSION}\label{discussion}

With the basic Dryad infrastructure in place, we see a number of
interesting future research directions. One fundamental question is the
applicability of the programming model to general large-scale
computations beyond text processing and relational queries. Of course,
not all programs are easily expressed using a coarse-grain data-parallel
communication graph, but we are now well positioned to identify and
evaluate Dryad's suitability for those that are.

Section 3 assumes an application developer will first construct a static
job graph, then pass it to the runtime to be executed. Section 6.3 shows
the benefits of allowing applications to perform automatic dynamic
refinement of the graph. We plan to extend this idea and also introduce
interfaces to simplify dynamic modifications of the graph according to
application-level control flow decisions. We are particularly interested
in data-dependent optimizations that might pick entirely different
strategies (for example choosing between in-memory and external sorts)
as the job progresses and the volume of data at intermediate stages
becomes known. Many of these strategies are already described in the
parallel database literature, but Dryad gives us a flexible testbed for
exploring them at very large scale. At the same time, we must ensure
that any new optimizations can be targeted by higher-level languages on
top of Dryad, and we plan to implement comprehensive support for
relational queries as suggested in Section 7.3.

The job manager described in this paper assumes it has exclusive control
of all of the computers in the cluster, and this makes it difficult to
efficiently run more than one job at a time. We have completed
preliminary experiments with a new implementation that allows multiple
jobs to cooperate when executing concurrently. We have found that this
makes much more efficient use of the resources of a large cluster, but
are still exploring variants of the basic design. A full analysis of our
experiments will be presented in a future publication.

There are many opportunities for improved performance monitoring and
debugging. Each run of a large Dryad job generates statistics on the
resource usage of thousands of executions of the same program on
different input data. These statistics are already used to detect and
re-execute slow-running ``outlier'' vertices. We plan to keep and
analyze the statistics from a large number of jobs to look for patterns
that can be used to predict the resource needs of vertices before they
are executed. By feeding these predictions to our scheduler, we may be
able to continue to make more efficient use of a shared cluster.

Much of the simplicity of the Dryad scheduler and fault-tolerance model
come from the assumption that vertices are deterministic. If an
application contains non-deterministic vertices then we might in future
aim for the guarantee that every terminating execution produces an
output that some failure-free execution could have generated. In the
general case where vertices can produce side-effects this might be very
hard to ensure automatically.

The Dryad system implements a general-purpose data-parallel execution
engine. We have demonstrated excellent scaling behavior on small
clusters, with absolute performance superior to a commercial database
system for a hand-coded read-only query. On a larger cluster we have
executed jobs containing hundreds of thousands of vertices, processing
many terabytes of input data in minutes, and we can automatically adapt
the computation to exploit network locality. We let developers easily
create large-scale distributed applications without requiring them to
master any concurrency techniques beyond being able to draw a graph of
the data dependencies of their algorithms. We sacrifice some
architectural simplicity compared with the MapReduce system design, but
in exchange we release developers from the burden of expressing their
code as a strict sequence of map, sort and reduce steps. We also allow
the programmer the freedom to specify the communication transport which,
for suitable tasks, delivers substantial performance gains.

\section{Acknowledgements}\label{acknowledgements}

We would like to thank all the members of the Cosmos team in Windows
Live Search for their support and collaboration, and particularly Sam
McKelvie for many helpful design discussions. Thanks to Jim Gray and our
anonymous reviewers for suggestions on improving the presentation of the
paper.

\section{10. REFERENCES}\label{references}

{[}1{]} Global grid forum. http://www.gridforum.org/.\\
{[}2{]} Intel IXP2XXX product line of network processors.
http://www.intel.com/design/network/products/npfamily/ixp2xxx.htm.\\
{[}3{]} Intel platform 2015.
http://www.intel.com/technology/architecture/platform2015/.

{[}4{]} The LINQ project.
http://msdn.microsoft.com/netframework/future/linq/.\\
{[}5{]} Open MPI. http://www.open-mpi.org/.\\
{[}6{]} SQL Server Integration Services.
http://www.microsoft.com/sql/technologies/integration/default.mspx.\\
{[}7{]} Thomas E. Anderson, David E. Culler, David A. Patterson, and NOW
Team. A case for networks of workstations: NOW. IEEE Micro, pages
54--64, February 1995.\\
{[}8{]} Remzi H. Arpaci-Dusseau. Run-time adaptation in River.
Transactions on Computer Systems (TOCS), 21(1):36--86, 2003.\\
{[}9{]} Özalp Babaoğlu, Lorenzo Alvisi, Alessandro Amoroso, Renzo
Davoli, and Luigi Alberto Giachini. Paralex: an environment for parallel
programming in distributed systems. pages 178--187, New York, NY, USA,
1992. ACM Press.\\
{[}10{]} Magdalena Balazinska, Hari Balakrishnan, Samuel Madden, and
Mike Stonebraker. Fault-Tolerance in the Borealis Distributed Stream
Processing System. In ACM SIGMOD, Baltimore, MD, June 2005.\\
{[}11{]} Tom Barclay, Robert Barnes, Jim Gray, and Prakash Sundaresan.
Loading databases using dataflow parallelism. SIGMOD Rec., 23(4):72--83,
1994.\\
{[}12{]} Chaitanya Baru and Gilles Fecteau. An overview of DB2 parallel
edition. In SIGMOD '95: Proceedings of the 1995 ACM SIGMOD international
conference on Management of data, pages 460--462, New York, NY, USA,
1995. ACM Press.\\
{[}13{]} Guy E. Blelloch. Programming parallel algorithms.
Communications of the ACM (CACM), 39(3):85--97, 1996.\\
{[}14{]} Robert D. Blumofe, Christopher F. Joerg, Bradley Kuszmaul,
Charles E. Leiserson, Keith H. Randall, and Yuli Zhou. Cilk: An
efficient multithreaded runtime system. In ACM SIGPLAN Symposium on
Principles and Practice of Parallel Programming (PPoPP), pages 207--216,
Santa Barbara, California, July 19-21 1995.\\
{[}15{]} Eylon Caspi, Michael Chu, Randy Huang, Joseph Yeh, Yury
Markovskiy, André DeHon, and John Wawrzynek. Stream computations
organized for reconfigurable execution (SCORE): Introduction and
tutorial. In FPL, Lecture Notes in Computer Science. Springer Verlag,
2000.\\
{[}16{]} Jeff Dean and Sanjay Ghemawat. MapReduce: Simplified data
processing on large clusters. In Proceedings of the 6th Symposium on
Operating Systems Design and Implementation (OSDI), pages 137--150,
December 2004.\\
{[}17{]} D. DeWitt, S. Ghandeharizadeh, D. Schneider, H. Hsiao, A.
Bricker, and R. Rasmussen. The GAMMA database machine project. IEEE
Transactions on Knowledge and Data Engineering, 2(1), 1990.\\
{[}18{]} David DeWitt and Jim Gray. Parallel database systems: The
future of high performance database processing. Communications of the
ACM, 36(6), 1992.\\
{[}19{]} D. C. DiNucci and R. G. Babb II. Design and implementation of
parallel programs with LGDF2. In Digest of Papers from Compcon '89,
pages 102--107, 1989.\\
{[}20{]} Armando Fox, Steven D. Gribble, Yatin Chawathe, Eric A. Brewer,
and Paul Gauthier. Cluster-based scalable network services. In ACM
Symposium on Operating Systems Principles (SOSP), pages 78--91, New
York, NY, USA, 1997. ACM Press.\\
{[}21{]} Sanjay Ghemawat, Howard Gobioff, and Shun-Tak Leung. The Google
file system. In SOSP '03: Proceedings of the nineteenth ACM symposium on

Operating systems principles, pages 29--43, New York, NY, USA, 2003. ACM
Press.\\
{[}22{]} Goetz Graefe. Encapsulation of parallelism in the Volcano query
processing system. In SIGMOD '90: Proceedings of the 1990 ACM SIGMOD
international conference on Management of data, pages 102--111, New
York, NY, USA, 1990. ACM Press.\\
{[}23{]} J. Gray, A.S. Szalay, A. Thakar, P. Kunszt, C. Stoughton, D.
Slutz, and J. Vandenberg. Data mining the SDSS SkyServer database. In
Distributed Data and Structures 4: Records of the 4th International
Meeting, pages 189--210, Paris, France, March 2002. Carleton Scientific.
also as MSR-TR-2002-01.\\
{[}24{]} Jim Gray and Alex Szalay. Science in an exponential world.
Nature, 440(23), March 23 2006.\\
{[}25{]} J.-H. Hwang, M. Balazinska, A. Rasin, U. Çetintemel, M.
Stonebraker, and S. Zdonik. A comparison of stream-oriented
high-availability algorithms. Technical Report TR-03-17, Computer
Science Department, Brown University, September 2003.\\
{[}26{]} Ujval Kapasi, William J. Dally, Scott Rixner, John D. Owens,
and Brucek Khailany. The Imagine stream processor. In Proceedings 2002
IEEE International Conference on Computer Design, pages 282--288,
September 2002.\\
{[}27{]} Eddie Kohler, Robert Morris, Benjie Chen, John Jannotti, and M.
Frans Kaashoek. The Click modular router. ACM Transactions on Computer
Systems, 18(3):263--297, 2000.\\
{[}28{]} James Larus and Michael Parkes. Using cohort scheduling to
enhance server performance. In Usenix Annual Technical Conference, June
2002.\\
{[}29{]} Orlando Loques, Julius Leite, and Enrique Vinicio Carrera E.
P-RIO: A modular parallel-programming environment. IEEE Concurrency,
6(1):47--57, 1998.\\
{[}30{]} William Mark, Steven Glanville, Kurt Akeley, and Mark J.
Kilgard. Cg: A system for programming graphics hardware in a C-like
language. ACM Transactions on Graphics, 22(3):896--907, 2003.\\
{[}31{]} P. Newton and J.C. Browne. The CODE 2.0 graphical parallel
programming language. pages 167 - 177, Washington, D. C., United States,
July 1992.\\
{[}32{]} Ken Phillips. SenSage ESA. SC Magazine, March 1 2006.\\
{[}33{]} Rob Pike, Sean Dorward, Robert Griesemer, and Sean Quinlan.
Interpreting the data: Parallel analysis with Sawzall. Scientific
Programming, 13(4):277--298, 2005.\\
{[}34{]} Mehul A. Shah, Joseph M. Hellerstein, and Eric Brewer. Highly
available, fault-tolerant, parallel dataflows. In SIGMOD '04:
Proceedings of the 2004 ACM SIGMOD international conference on
Management of data, pages 827--838, New York, NY, USA, 2004. ACM
Press.\\
{[}35{]} V. S. Sunderam. PVM: a framework for parallel distributed
computing. Concurrency: Pract. Exper., 2(4):315--339, 1990.\\
{[}36{]} David Tarditi, Sidd Puri, and Jose Oglesby. Accelerator: using
data-parallelism to program GPUs for general-purpose uses. In
International Conference on Architectural Support for Programming
Languages and Operating Systems (ASPLOS), Boston, MA, October 2006. also
as MSR-TR-2005-184.\\
{[}37{]} Douglas Thain, Todd Tannenbaum, and Miron Livny. Distributed
computing in practice: The Condor experience. Concurrency and
Computation: Practice and Experience, 17(2-4):323--356, 2005.\\
{[}38{]} P.W. Trinder, H-W. Loidl, and R.F. Pointon. Parallel and
distributed Haskells. Journal of Functional Programming,
12(4\&5):469--510, 2002.

\end{document}
